\documentclass[11pt]{ctexart}
\usepackage[margin=1.12in]{geometry}
\usepackage{amsmath,amssymb,amsthm,mathtools,mathrsfs}
\usepackage{microtype}
\usepackage{enumitem}
\usepackage{xcolor}
\usepackage{tikz}
\usepackage{hyperref}
\hypersetup{colorlinks=true,linkcolor=blue!45!black,citecolor=blue!45!black,urlcolor=blue!45!black,pdftitle={最优 Max-Cut 难度的直接证明},pdfauthor={OpenAI（中文对照版）}}
\usepackage[nameinlink,capitalise,noabbrev]{cleveref}
% —— 中文环境名 ——
\newtheorem{theorem}{定理}[section]
\newtheorem{lemma}[theorem]{引理}
\newtheorem{proposition}[theorem]{命题}
\newtheorem{corollary}[theorem]{推论}
\theoremstyle{definition}
\newtheorem{definition}[theorem]{定义}
\theoremstyle{remark}
\newtheorem{remark}[theorem]{注记}
\crefname{theorem}{定理}{定理}
\crefname{lemma}{引理}{引理}
\crefname{proposition}{命题}{命题}
\crefname{corollary}{推论}{推论}
\crefname{definition}{定义}{定义}
\crefname{remark}{注记}{注记}
\crefname{equation}{式}{式}
\crefname{section}{节}{节}
\crefname{appendix}{附录}{附录}
\crefname{figure}{图}{图}
\crefname{table}{表}{表}
\numberwithin{equation}{section}
\newcommand{\F}{\mathbb F_2}
\newcommand{\E}{\mathbb E}
\newcommand{\one}{\mathbf 1}
\newcommand{\GW}{\alpha_{\mathrm{GW}}}
\DeclareMathOperator{\Aff}{Aff}
\DeclareMathOperator{\Stab}{Stab}
\DeclareMathOperator{\Inf}{Inf}
\DeclareMathOperator{\Val}{Val}
\DeclareMathOperator{\TV}{TV}
\DeclareMathOperator{\supp}{supp}
\DeclareMathOperator{\codim}{codim}
\DeclarePairedDelimiter{\abs}{\lvert}{\rvert}
\DeclarePairedDelimiter{\norm}{\lVert}{\rVert}
\newcommand{\ip}[2]{\langle #1,#2\rangle}
\newcommand{\bits}{\{-1,1\}}
\newcommand{\Bfun}{B}
\newcommand{\cut}{\operatorname{cut}}
\newcommand{\eps}{\varepsilon}
\newcommand{\ind}{\mathbf 1}
\setlist{topsep=4pt,itemsep=2pt}
\allowdisplaybreaks[2]

\title{最优 Max-Cut 难度的直接证明\\{\large A Direct Proof of Optimal Max-Cut Hardness（中文对照版）}}
\author{OpenAI（原文）\\中文翻译：openai-math-zh 项目}
\date{2026 年 9 月 23 日（原文日期）}

\begin{document}
\maketitle
\begin{abstract}
我们证明：在简单无权图上，把 Max-Cut 近似到任何优于
Goemans--Williamson 常数的固定因子都是 NP 难的。

\smallskip
\emph{（本文件是 openai/math 手稿的中文对照翻译；原文结果未经同行评审。
数学公式、定理编号与结构与英文原文保持一致，术语以仓库 glossary.md 为准；
译文与原文冲突时以英文原文为准。）}
\end{abstract}

\section{引言}

对于边权非负的图，Max-Cut（最大割）问的是：如何划分顶点，使横跨划分的边权之和最大。
一个割由顶点上的符号编码，一条边恰在其两端符号不同时有贡献。
对最大化问题而言，$\alpha$-近似指返回一个价值至少为最优值 $\alpha$ 倍的可行解。
除非特别说明，图及其有理边权都是显式给出的。
我们以 $\operatorname{MaxCut}(G)$ 记最大跨越权重；当边权之和为 1 时，
也把该值记作 $\Val(G)$。
带权图的割是否至少达到给定权重的判定问题，出现在 Karp 最初的
NP 完全问题清单中~\cite{Karp72}。近似算法问的是：在多项式时间内，
这个最优值还能保证多少。
Goemans 与 Williamson~\cite{GW95} 为 Max-Cut 给出了半定规划算法与
随机超平面舍入，其期望近似比逼近
\begin{equation}\label{eq:gw-constant}
 \GW=\min_{-1\le \rho<1}
 \frac{2\arccos\rho}{\pi(1-\rho)}
 =0.878567\ldots.
\end{equation}
该表达式比较的是端点向量内积为 $\rho$ 的一条边的两项贡献：其在半定目标中
的价值为 $(1-\rho)/2$，而随机超平面以概率 $\arccos(\rho)/\pi$ 分开这两个
向量。比值对 $\rho$ 取最小，就是这一舍入规则的一致保证。
在普通 NP 难度意义下判定该比值是否最优，一直是近似算法领域的核心问题。
H\aa stad 的难度定理~\cite{Hastad01} 结合 Trevisan、Sorkin、Sudan 与
Williamson 的归约 gadget~\cite{TSSW00}，排除了大于 $16/17$ 的比值。
Khot、Kindler、Mossel 与 O'Donnell~\cite{KKMO07} 在唯一游戏
猜想~\cite{Khot02} 的假设下确立了 $\GW$ 的最优性；其解析输入是
Mossel、O'Donnell 与 Oleszkiewicz 的 Majority Is Stablest
定理~\cite{MOO10}。

我们解决 Max-Cut 无条件最优性的问题。

\begin{theorem}\label{thm:main}
对每个满足 $\GW<\alpha\le1$ 的固定 $\alpha$，在因子 $\alpha$ 内近似
Max-Cut 是 NP 难的。该结论对简单无权图成立。
\end{theorem}

这一陈述的含义是：若存在多项式时间算法总能返回权重至少为最大值
$\alpha$ 倍的割，则 $\mathrm P=\mathrm{NP}$。我们证明下面更精确的
间隙定理。记
\begin{equation}\label{eq:B}
 B(t)=\frac{2}{\pi}\arcsin t,\qquad 0<t<1.
\end{equation}

\begin{theorem}\label{thm:gap}
固定有理数 $t\in(0,1)$ 与常数 $0<\eps<(t-B(t))/4$。对边上有有理概率
分布的图，区分下列情形是 NP 难的：
\begin{align*}
 \textnormal{YES：}\quad&\Val\ge\frac{1+t}{2}-\eps,\\
 \textnormal{NO：}\quad&\Val\le\frac{1+B(t)}{2}+\eps,
\end{align*}
其中 $\Val$ 为最大割权重。
\end{theorem}

\subsection*{已有工作与直接路线}

把证明验证归约为近似间隙的思想可追溯到 Feige、Goldwasser、Lov\'asz、
Safra 与 Szegedy~\cite{FGLSS96}，以及 Arora--Safra 与 Arora、Lund、
Motwani、Sudan、Szegedy 的 PCP 定理~\cite{AroraSafra98,ALMSS98}。
H\aa stad 对验证测试的傅里叶分析~\cite{Hastad01} 使其中许多间隙达到
锐利。对 Max-Cut 而言，半定松弛的高斯几何给出了一个格外精确的目标。
Feige 与 Schechtman~\cite{FS02} 为该松弛构造了逼近 $\GW$ 的整性间隙；
但仅凭该松弛的局限，并不能对任意算法建立难度。
Khot、Kindler、Mossel 与 O'Donnell 的唯一游戏归约~\cite{KKMO07} 把同一
几何阈值与计算难度联系起来。

Khot、Kindler、Mossel 与 O'Donnell 的原始分析提出：要得到该阈值，
所需的解析命题是 Majority Is Stablest。Mossel、O'Donnell 与
Oleszkiewicz 通过在低影响布尔函数与高斯函数之间作比较证明了它
~\cite{MOO10}，把 Borell 的半空间噪声稳定性界~\cite{Borell85} 转移到
布尔情形。他们的定理使唯一游戏猜想仍是经典难度归约中的假设
~\cite[Corollary~2.12]{MOO10}。我们用 Majority Is Stablest 找出测试中
的一个具有影响的坐标，再证明所需的额外抽取步骤，把该坐标变成
源博弈（source game）的标记。上述算法基准是理想舍入比；有限精度的
SDP 优化与随机超平面舍入对任意固定 $\eta>0$ 给出期望比
$\GW-\eta$~\cite[Section~2]{GW95}。

半定近似与难度之间的比较不止于单一的最坏情形比值。O'Donnell 与
Wu~\cite{OW08} 在唯一游戏猜想下确定了 Max-Cut 的完整近似曲线，把可
达到的割值与实例的最优值联系起来。Raghavendra~\cite{Raghavendra08}
对有限约束满足问题建立了半定松弛与基于唯一游戏的难度之间的一般对
应。本定理关心的是 Max-Cut 无条件的最坏情形比值；下文的证明建立了
\Cref{thm:gap} 中特定的完备性与可靠性间隙族。

在投影博弈（projection game）中，二部图的每个顶点从指定字母表中领取
标记，边权构成一个概率分布。当边所规定的映射把左端标记送到右端标记
时，该边被接受。其值是被接受边权的最大值。二对一约束指每个右侧标记
恰有两个原像。完备性与可靠性分别是 YES 情形值的承诺下界与 NO 情形值
的上界。

我们的出发点是已发表的、完备性不完美的二对一博弈构造~\cite{DKKMS25}。
Grassmann 图方法由 Khot、Minzer 与 Safra~\cite{KMS25} 引入。其可靠性
经由 Barak、Kothari 与 Steurer 建立的扩张与一致性之间的联系~\cite{BKS19}
以及 Khot、Minzer 与 Safra 的扩张定理~\cite{KMS23} 得以补全。我们
需要该构造的一个精确仿射形式：其投影是维数相差 1 的二元空间之间的
满射仿射映射，且在把完备性误差取得任意小之前，字母表可以先按所需
可靠性固定。\Cref{prop:affine-games} 给出该输入。
Appendix~\ref{sec:affine-games} 从已发表构造推导出它，同时指定
有效子空间的条件化与可靠性证明中使用的辅助候选列表。
与 Majority Is Stablest 一起，该构造构成了本论证的外部基础。
唯一游戏猜想只作为历史背景出现。树的傅里叶分析、哈希界与隐藏坐标
解码均在此处证明。

\subsection*{证明及其主要组成部分}

如同长码（long code）归约~\cite{BGS98,Hastad01}，标记 $a$ 由独裁者
函数 $w\mapsto w(a)$ 表示，其中 $w$ 取遍以标记为指标的布尔字。
图的每个左顶点元组对应一个这样的字域。于是任意割在每个域上指定一个
布尔函数。为采样一条边，测试通过共同的右元组选取两个左元组，在右侧
标记空间上生成两个字，再沿采样到的仿射投影拉回；然后把第二个字取反。
当采样到的约束被满足时，独裁者标记给出的割概率为二分之一加一：
即两字在公共标记处取值相关再加二分之一。

字生成器是一棵有限有根树。每个内部节点有 $s$ 对子节点以及一张从
$\F^s$ 到 $\bits$ 的随机表；叶子承载标记空间上的仿射型。在标记 $a$
处，这些仿射型给出叶符号，各表递归地把子符号组合成根符号 $W(a)$。
门规则选取的原则是：其 Walsh 展开的每个单项式从每对子节点中恰好保留
一个。第二次探索在每个到达的节点随机选取一对并保留其两个子节点。
在任何固定深度，两次探索恰有一个公共节点。测试把第二次探索在随机
深度处到达的那些表耦合起来，逐项相关取为 $t$，其余所有表保持不变。
于是每个单项式恰获得一个因子 $t$。对仿射行做一次小规模的独立重采样
提供了可靠性所需的平滑，且只带来很小的完备性损失。

对任意割，我们先把它在源博弈边上的奇部取平均。所得函数在选定的表
项上是均衡的；若割高于所需的 NO 阈值，则其期望噪声稳定性超过
$B(t)$。于是 Majority Is Stablest 检出一个有影响的项。该项由门的输入
索引，而不是由源博弈标记索引。剩下的任务是：把该影响与源边的两个
端点能够一致选取的标记连接起来。分三步完成。

\begin{enumerate}
\item \emph{傅里叶支撑的局部化简。}
门的对称性迫使每对中的两个子指标具有相同的和。在行平滑之后，对选定
表的深度取平均，使得只保留「把该和放进每对的一个子节点、在另一个中
置零」的指标变得不再昂贵。这一化简只在有影响的门处施加。正是这一
结构让仿射平移表现为该门的表的平移。

\item \emph{向量值仿射哈希引理。}
对独立的仿射平移，弥散分布的傅里叶质量无法在平均意义下产生大的单项
影响。该界对行空间的维数一致。其证明使用两次截断与一个四阶矩恒等式；
凡需在不损失字母表规模的前提下解码平移系数阵列的场合，它都可能有用。

\item \emph{带隐藏坐标的解码。}
质量集中是在固定许多行坐标之后得到的。左端点通常无法在不知道其关联
投影的情况下构造该上下文的拉回。我们使用博弈的乘积，并在上下文与
潜在重采样行的某个均匀选取的因子中植入零线性系数。在该因子被遗忘后，
植入律接近原律。在植入因子处，上下文可以在不知道投影的情况下拉回。
只有选定门处的 $2s$ 个自由行仍需跨该因子比较，其代价是傅里叶质量
损失 $2^{2s}$。
\end{enumerate}

由此得到的解码损失与源字母表无关。这一点至关重要：字母表维数依赖
于我们要求的可靠性，若损失依赖维数，可靠性的选取就会循环。我们先
固定该损失与源可靠性，再固定字母表与乘积因子个数，最后固定源完备性
误差。\Cref{prop:affine-games} 中的量词顺序恰好允许这一序列。

\Cref{sec:preliminaries} 固定傅里叶约定。
\Cref{sec:tree,sec:reduction} 定义编码与 Max-Cut 测试。
\Cref{sec:tree-fourier,sec:hashing,sec:extraction} 证明解析组成部分。
\Cref{sec:decoding,sec:completion} 给出解码并完成证明。
Appendix~\ref{sec:affine-games} 从已发表构造推导精确的起始博弈，
Appendix~\ref{sec:unweighted} 给出到简单无权图的确定性转换。

\section{预备知识}\label{sec:preliminaries}

证明使用两种傅里叶展开。仿射行的特征刻画候选标记，而布尔表项的特征
检测 Majority Is Stablest 所适用的那些影响。我们先介绍这两套约定，
再陈述精确的仿射博弈输入。

$\F$ 上所有向量空间都是有限维的。除非另有说明，有限向量空间或布尔
立方上的所有期望都使用均匀概率测度。所有傅里叶系数都使用这一概率
归一化。$L^2$ 范数始终对其下标显示的变量积分；声明为固定的变量不
参与积分。

\subsection{仿射行与标记}

设 $\mathcal A$ 是 $\F$ 上的非空仿射空间，带固定的仿射坐标系，并令
\[
 H=\Aff(\mathcal A,\F).
\]
这是一个向量空间，包含常值型 $\one$。
对 $q\in H^*$ 定义 $\chi_q(h)=(-1)^{q(h)}$。
求值把 $\mathcal A$ 等同于
\begin{equation}\label{eq:labels-dual}
 \{q\in H^*:q(\one)=1\}.
\end{equation}
事实上，若 $h(a)=h_0+\sum_i h_i a_i$，则在 $\one$ 上取值 1 的线性
泛函具有唯一形式 $q(h)=h_0+\sum_i h_i a_i$。
我们恰把这些泛函称为\emph{标记}；$H^*$ 中的一般非零泛函未必是标记。

对 $f:H^k\to\mathbb R$ 与 $C=(q_1,\ldots,q_k)\in(H^*)^k$，令
\[
 \widehat f(C)=\E_R f(R)\prod_{i=1}^k\chi_{q_i}(R_i),\qquad
 |C|=\#\{i:q_i\ne0\}.
\]
Parseval 给出 $\sum_C\widehat f(C)^2=\E_R f(R)^2$。
带额外变量的函数沿用同一记号；这些变量在所指的傅里叶展开中保持
固定。

对 $0\le r\le1$，记 $T_r$ 为行噪声算子：在每个坐标独立地以概率
$r$ 保留该行，否则代之以独立的均匀行。则
\begin{equation}\label{eq:row-noise}
 \widehat{T_rf}(C)=r^{|C|}\widehat f(C).
\end{equation}
该算子保持值域 $[-1,1]$。它也是与独立增量 $E\in H^k$ 的卷积，
其坐标以概率 $r$ 为零、以概率 $1-r$ 为独立均匀行。
卷积与行的任何固定平移交换。

\subsection{布尔傅里叶分析}

对有限集 $\Omega$ 与 $h:\bits^\Omega\to\mathbb R$，写
\[
 h(x)=\sum_{S\subseteq\Omega}\widehat h(S)\prod_{a\in S}x_a.
\]
其在 $a\in\Omega$ 处的影响与次数至多 $K$ 的影响为
\[
 \Inf_a(h)=\sum_{S\ni a}\widehat h(S)^2,\qquad
 \Inf_a^{\le K}(h)=\sum_{\substack{S\ni a\\|S|\le K}}
 \widehat h(S)^2.
\]
记 $P_{\le K}$ 为到傅里叶次数至多 $K$ 的正交投影。特别地，
\begin{equation}\label{eq:degree-influence}
 \sum_{a\in\Omega}\Inf_a(P_{\le K}h)
 \le K\norm{h}_2^2.
\end{equation}
对 $t$-相关的均匀比特 $x,y$（在每个坐标独立耦合），
\begin{equation}\label{eq:stability}
 \Stab_t(h)=\E[h(x)h(y)]
 =\sum_{S\subseteq\Omega}t^{|S|}\widehat h(S)^2.
\end{equation}
因此当 $t\ge0$ 时稳定性非负。若 $h\in[-1,1]$，则它至多为 1。

\begin{lemma}[带次数截断的 Majority Is Stablest]\label{lem:mis}
固定 $0<t<1$ 与 $\xi>0$。存在只依赖 $t,\xi$ 的整数 $K\ge1$ 与
$\tau>0$，使得每个满足 $\E h=0$ 的 $h:\bits^\Omega\to[-1,1]$ 都有
\[
 \bigl(\Inf_a^{\le K}(h)<\tau\text{ 对所有 }a\in\Omega\bigr)
 \quad\Longrightarrow\quad
 \Stab_t(h)\le B(t)+\xi.
\]
\end{lemma}

\begin{proof}
无截断的断言来自 Majority Is Stablest~\cite[Theorem~4.4]{MOO10}。
为匹配其 $[0,1]$ 约定，令 $g=(h+1)/2$：则 $\E g=1/2$，每个非常数
傅里叶系数减半，且 $\Stab_t(h)=4\Stab_t(g)-1$。$g$ 的高斯半空间值
为 $1/4+\arcsin(t)/(2\pi)$，在我们的归一化下即为 $B(t)$；见
\cite[Theorem~4.1 and Corollary~4.3]{MOO10} 与 \cite{Borell85}。
下面是从无截断断言到带次数截断版本的归约。对立方噪声参数 $u<1$，
令 $h_u=T_u h$（本段使用普通比特噪声）。值域与均值保持不变。对范数
至多 1 的 $h$ 一致地有
\[
 0\le\Stab_t(h)-\Stab_t(h_u)
 \le\sup_{j\ge0}t^j(1-u^{2j})\longrightarrow0
 \quad(u\uparrow1).
\]
收敛性这样得出：先选有限的次数截断以控制尾部 $t^j$，再对剩余次数取
$u$ 接近 1。选 $u$ 使该误差至多 $\xi/2$。设 $\delta>0$ 为误差
$\xi/2$ 对应的无截断影响阈值。选 $K$ 使 $u^{2(K+1)}\le\delta/2$，
并取 $\tau=\delta/2$。则
\[
 \Inf_a(h_u)
 \le\Inf_a^{\le K}(h)+u^{2(K+1)}\norm{h}_2^2
 <\delta.
\]
对 $h_u$ 应用无截断定理并加上稳定性误差即得。
\end{proof}

\subsection{仿射投影博弈}

带权二部投影博弈由以下部分组成：有限集 $X,Y$；有限边集 $E$ 上的
有理概率分布 $\lambda$，边带指向 $X$ 与 $Y$ 的端点映射（允许平行
边）；有限字母表；以及每条边上的投影 $\pi_e$。其值为
\[
 \Val(\mathcal G)=
 \max_{\ell_X,\ell_Y}\Pr_{e=(x,y)\sim\lambda}
 [\pi_e(\ell_X(x))=\ell_Y(y)].
\]
本文中字母表为 $\F^{n+1}$ 与 $\F^n$，且每个 $\pi_e:\F^{n+1}\to\F^n$
都是满射仿射映射。\Cref{prop:affine-games} 提供具有该结构的 NP 难
间隙。凡使用条件边分布之处，我们一律弃掉边际概率为零的端点。

起始定理的量词顺序为
\[
 \forall b\in(0,1)\quad \exists n=n(b)\ge1\quad
 \forall a\in(0,1-b).
\]
精确承诺见下；其从已发表构造的推导在 \Cref{sec:affine-games}。

\begin{proposition}[仿射二对一博弈定理]
\label{prop:affine-games}
对每个常数 $b\in(0,1)$ 存在整数 $n\geq 1$，使得对每个常数
$a\in(0,1-b)$，下述承诺问题是 NP 难的。其输入是带权二部博弈：顶点集
$X,Y$，各自的标记集 $\F^{n+1},\F^n$，以及每条边 $e=(x,y)$ 上的满射
仿射映射
\[
  \pi_e:\F^{n+1}\longrightarrow\F^n.
\]
当 $\pi_e(\lambda_X(x))=\lambda_Y(y)$ 时称该边被满足。承诺情形为
\[
  \mathrm{YES}：\quad \Val(\mathcal G)\geq 1-a,
  \qquad
  \mathrm{NO}：\quad \Val(\mathcal G)\leq b.
\]
整数 $n$ 只依赖 $b$。对固定的 $a,b$，归约是确定性的，运行时间与
输出规模均为多项式。
\end{proposition}

带权图同样由无序边（允许自环）上的有理分布给出。割给顶点指派符号，
当边两端符号不同时该边被满足；其值即被满足概率。自环的割权重为零。
\Cref{lem:unweighted} 中的转换即使在删除自环后也保持公共尺度。

\section{带互补遍历的树码}
\label{sec:tree}

我们沿用 \Cref{sec:preliminaries} 的仿射标记空间 $\mathcal A$ 与行空间
$H=\Aff(\mathcal A,\F)$。
下述编码把仿射行与布尔门表变成一个以标记为指标的词。它的两种互补
遍历控制了在单一标记处取词值的割所看到的相关性。

\begin{definition}[树码]
\label{def:tree-code}
固定整数 $s\ge2$ 与 $m\ge1$，并令
\[
 k=(2s)^m,\qquad d=s^m.
\]
设 $\mathcal T$ 是深度为 $m$ 的完全有根 $2s$ 叉树。
每个内部节点 $v$ 的子节点由 $(g,b)\in[s]\times\{0,1\}$ 索引。其叶子
由 $[k]$ 索引。给每个内部节点指派一张表 $F_v:\F^s\to\{-1,1\}$，
给每个叶子指派一行 $R_i\in H$。
词 $W(R,F)\in\{-1,1\}^{\mathcal A}$ 逐点定义如下。在 $a\in\mathcal A$
处，叶子 $i$ 的符号为 $(-1)^{R_i(a)}$。子符号为 $(Z_{g,b})$ 的内部
节点输出
\begin{equation}
 \left(\prod_{g=1}^s Z_{g,0}\right)F_v(z),
 \qquad (-1)^{z_g}=Z_{g,0}Z_{g,1}.
 \label{eq:gate-rule}
\end{equation}
根节点的输出即 $W(R,F)(a)$。
\end{definition}

我们让这些表被树的第二次遍历所噪声化。
\emph{随机切片}的取法是：先从 $\{0,\ldots,m-1\}$ 中均匀选取 $J_0$。
从根出发，在深度小于 $J_0$ 的每个到达节点处，均匀且独立地选取其
$s$ 对子节点中的一对，并沿该对的两个子节点继续。在深度 $J_0$ 处
到达的节点集 $\mathcal D$ 即为切片。

在 $J_0=j$ 的条件下，切片有 $2^j$ 个节点，而每个固定的深度 $j$ 节点
以概率 $s^{-j}$ 属于它。记 $F_{\rm out}=(F_u)_{u\notin\mathcal D}$ 为
切片外的表。

下文用到的树的两种遍历互为补充。切片构造选取一对并保留两个子节点；
门的 Walsh 展开从每对中保留一个子节点。因此二者的交沿一条单一路径。
Figure~\ref{fig:traversals} 展示了局部规则；逐层迭代即在每个切片深度
得到一个公共节点。

\begin{figure}[ht]
\centering
\begin{tikzpicture}[x=1cm,y=1cm,every node/.style={font=\small},
  mono/.style={draw=orange!80!black,line width=1.4pt},
  slice/.style={draw=blue!65!black,line width=1.4pt,dashed}]
\node[circle,draw,minimum size=7mm] (v) at (0,1.7) {$v$};
\foreach \x/\name/\lab in {-3/a/{(1,0)},-2/b/{(1,1)},-.5/c/{(2,0)},.5/d/{(2,1)},2/e/{(3,0)},3/f/{(3,1)}} {
 \node[circle,draw,fill=white,minimum size=3mm,inner sep=1pt] (\name) at (\x,0) {};
 \node[below] at (\x,-.12) {$\lab$};
 \draw[gray!40] (v) -- (\name);
}
\draw[mono] (v) -- (a);
\draw[mono] (v) -- (d);
\draw[mono] (v) -- (e);
\draw[slice] (v) -- (c);
\draw[slice] (v) -- (d);
\node[circle,draw,fill=violet!65,minimum size=4mm,inner sep=1pt] at (.5,0) {};
\draw[mono] (-3,-1) -- (-2.4,-1);
\node[anchor=west] at (-2.3,-1) {每对一个子节点};
\draw[slice] (-3,-1.5) -- (-2.4,-1.5);
\node[anchor=west] at (-2.3,-1.5) {一对的两个子节点};
\end{tikzpicture}

\caption{一个门处的互补选择，以 $s=3$ 为例。单项式遍历（实线）从每对
中选一个子节点；切片探索（虚线）选一对的两个子节点。恰有一个子节点
被共享，以紫色标出。图中只描绘这一局部选择；同一规则在所选切片之上
逐层重复。}
\label{fig:traversals}
\end{figure}

\begin{lemma}[互补遍历与相关性]
\label{lem:tree-correlation}
对每组固定的表 $F$，根输出作为 $k$ 个叶符号的多重线性多项式只含
次数为 $d$ 的单项式。每个这样的单项式在其到达的每个节点处都从每对
中选一个子节点。每条如此得到的遍历与每个可能的切片 $\mathcal D$
恰交于一个节点。因此，把某一切片上的所有表取反即把 $W$ 取反。

设 $0<t<1$ 与 $0<r<1$。均匀抽取 $F$，并如下构造 $F'$：切片外的表
保持不变，而 $\mathcal D$ 上每对相应的表项耦合为相关为 $t$ 的均匀
符号，各表项间独立。独立地均匀抽取 $R\in H^k$，并以概率 $r$ 保留
$R$ 的每行、否则代之以独立的均匀行（在两个副本中独立进行），得到
$R^1,R^2$。则对每个 $a\in\mathcal A$ 与每个固定的切片，
\begin{equation}
 \E\bigl[W(R^1,F)(a)W(R^2,F')(a)\mid\mathcal D\bigr]
       =t r^{2d}.
 \label{eq:tree-correlation}
\end{equation}
特别地，若 $r=1-\mu/d$（$0<\mu<1/2$），则 $r^{2d}\ge1-2\mu$。
\end{lemma}

\begin{proof}
对固定的门表，按自变量 $z$ 展开：
$F_v(z)=\sum_{I\subseteq[s]}\widehat F_v(I)(-1)^{\sum_{g\in I}z_g}$。
代入 Equation~\eqref{eq:gate-rule}，指标为 $I$ 的项是
\[
 \widehat F_v(I)
 \prod_{g\notin I}Z_{g,0}\prod_{g\in I}Z_{g,1}.
\]
迭代此展开给出
\begin{equation}
 W_F(x)=\sum_{P\in\mathscr P}a_P(F)\prod_{i\in P}x_i,
 \qquad |P|=d,
 \label{eq:tree-polynomial}
\end{equation}
其中 $\mathscr P$ 是在每个到达的内部节点处从每对中选一个子节点所得
叶子集的族。叶子集决定所有这些选择：被选的子节点与该叶子集之交
非空，未选的则不相交。因此不同的遍历给出不同的单项式，而
$a_P(F)$ 是遍历到达的每个节点处 $F_v$ 的一个 Walsh 系数之积。

在根处，单项式遍历与切片探索同时存在。在切片深度之上的任何公共
节点处，两种探索恰选同一个子节点。归纳可得在切片深度处恰有一个
公共节点。于是每个 $a_P$ 恰含来自 $\mathcal D$ 上某张表的一个因子。
把所有这样的表取反，就把 Equation~\eqref{eq:tree-polynomial} 中的
每个系数取反，从而把词取反。此外，切片表的每个 Walsh 系数满足
\[
 \E[\widehat F'_v(I)\mid F_v]=t\widehat F_v(I).
\]
各表之间的条件独立性给出
$\E[a_P(F')\mid F,\mathcal D]=t a_P(F)$。

对固定的 $a$，符号 $x_i^b=(-1)^{R_i^b(a)}$ 均匀且关于 $i$ 独立，
且 $\E[x_i^1x_i^2]=r^2$。于是
Equation~\eqref{eq:tree-polynomial} 中两个单项式乘积的期望在其叶子
集不一致时为零，一致时为 $r^{2d}$。由于 $W_F$ 是布尔值，Parseval
给出 $\sum_P a_P(F)^2=1$。先对行取平均、再对 $F'$ 取平均，即证
Equation~\eqref{eq:tree-correlation}。最后一断言是 Bernoulli
不等式：$(1-\mu/d)^{2d}\ge1-2\mu$。
\end{proof}

该相关性恒等式描述的是在单一标记处取词值的割。为分析任意割，我们
还需要能依赖整张词、却看不到其生成行与表的函数。若
\begin{equation}
 f(R,F)=\phi(W(R,F))
 \label{eq:function-of-code}
\end{equation}
对某个固定的词上函数 $\phi$ 成立，则称 $f$ 为\emph{码的函数}。
这一类在先固定 $\phi$ 再取平均下封闭，包括先把词经固定仿射映射拉回
的取法。对 $r\in[0,1]$，记 $T_r$ 为保留概率 $r$ 的独立行重采样，
并写 $\bar f=T_rf$。该算子保持 $f$ 的值域界。
若 $\phi$ 是奇函数，Lemma~\ref{lem:tree-correlation} 表明 $f$ 与
$\bar f$ 在切片表同时取反下都是奇的。特别地，当所有行与切片外的表
固定时，它们在这些表项上均值为零。

在这些固定数据下，其余随机输入是 $|\mathcal D|2^s$ 个独立表项构成
的立方，由 $(v,z)\in\mathcal D\times\F^s$ 索引。此立方上的傅里叶
次数与影响指\emph{表项}的子集。这与相关性证明中固定门在其 $s$ 比特
自变量 $z$ 上的 Walsh 展开不同。

\section{Max-Cut 归约}\label{sec:reduction}

固定有理数 $t\in(0,1)$、树参数 $s,m$、有理数 $0<\mu<1/2$ 与
$r=1-\mu/d$。这些参数将在 \Cref{sec:completion} 中选定。设
$\mathcal G$ 是 \Cref{prop:affine-games} 中的仿射博弈，边律为
$\lambda$，字母表维数为 $n+1,n$。设 $T\ge1$ 为整数。我们使用乘积
标记空间
\[
 \mathcal B=(\F^{n+1})^T,\qquad
 \mathcal A=(\F^n)^T,
\]
以及行空间 $H=\Aff(\mathcal A,\F)$。

\subsection{顶点与边分布}

对每个 $\underline x\in X^T$，为每个词 $w\in\bits^{\mathcal B}$ 建一个
图顶点 $(\underline x,w)$。这是标准的长码字域~\cite{BGS98,Hastad01}；
树构造规定测试如何采样其查询。于是任意割恰是一族函数
\[
 \sigma_{\underline x}:\bits^{\mathcal B}\longrightarrow\bits.
\]
这些函数只接收所得的词；它们不接收生成词的行、门表、切片或采样到的
博弈边。

图的边分布是下述测试。
\begin{enumerate}
\item 从 $\lambda$ 的乘积右端点边际中抽取 $\underline y\in Y^T$。
独立地由给定 $\underline y$ 的乘积条件边律抽取两条元组
$\underline e,\underline e'$。记其左端点为 $\underline x,\underline x'$，
其逐坐标投影为 $\pi_{\underline e},\pi_{\underline e'}:\mathcal B\to\mathcal A$。

\item 与所有博弈数据独立地抽取随机切片 $\mathcal D$。按
\Cref{lem:tree-correlation} 抽取 $F,F'$：在 $\mathcal D$ 上相应表项
为独立的 $t$-相关均匀对，在 $\mathcal D$ 外为相同的均匀对。

\item 抽取均匀的公共行数组 $R\in H^k$。在两个副本中各自独立地以概率
$r$ 保留每行、否则均匀重采样，得到 $R^1,R^2$。

\item 输出以
\begin{equation}\label{eq:test-queries}
 \bigl(\underline x,W(R^1,F)\circ\pi_{\underline e}\bigr),
 \qquad
 \bigl(\underline x',-W(R^2,F')\circ\pi_{\underline e'}\bigr)
\end{equation}
为端点的无序边。
\end{enumerate}
所有边权即为这些结果的概率；相同无序端点的结果被合并聚合。

\begin{lemma}[完备性]\label{lem:completeness}
若 $\Val(\mathcal G)\ge1-a$，则图有值至少为
\begin{equation}\label{eq:completeness}
 (1-2Ta)\frac{1+t r^{2d}}2
 \ge\frac{1+t}{2}-t\mu-2Ta
\end{equation}
的割。第一个下界在 $2Ta\le1$ 时有用，这正是下文使用的范围。
\end{lemma}

\begin{proof}
选取以概率至少 $1-a$ 满足边的博弈标记 $\ell_X,\ell_Y$。用左端标记
元组处的取值为每个词顶点染色：
\[
 \sigma_{\underline x}(w)=
 w\bigl(\ell_X(x_1),\ldots,\ell_X(x_T)\bigr).
\]
采样到的边元组的每个坐标都具有原边律。并界表明：所有 $2T$ 个采样
约束都以至少 $1-2Ta$ 的概率被满足。当所有约束被满足时，两个被查询
的元组投影到同一标记
$l=(\ell_Y(y_1),\ldots,\ell_Y(y_T))\in\mathcal A$。此时割符号之积为
$-W(R^1,F)(l)W(R^2,F')(l)$。词生成数据与博弈边独立，故
\Cref{lem:tree-correlation} 给出条件割概率 $(1+t r^{2d})/2$。在其余
事件上割概率非负。当 $2Ta\le1$ 时即得第一个界；若 $2Ta>1$，该界为
负，自动成立。最后，$r^{2d}\ge1-2\mu$，且前因子带来的损失至多
$2Ta$，得第二个界。
\end{proof}

\subsection{大割产生超额稳定性}

任意割在其被查询的词取反时未必变号。我们把它分解为两个奇偶部，因为
第二个查询前的负号使奇性依赖增加割值、偶性依赖减少割值。对每张割表
定义其偶部与奇部：
\[
 \sigma^{\mathrm{even}}(w)=\frac{\sigma(w)+\sigma(-w)}2,
 \qquad
 \sigma^o(w)=\frac{\sigma(w)-\sigma(-w)}2.
\]
对边元组与右端点元组，令
\begin{align}
 G_{\underline e}(R,F)
 &=\sigma^o_{\underline x(\underline e)}
       (W(R,F)\circ\pi_{\underline e}),\label{eq:edge-function}\\
 G_{\underline y}(R,F)
 &=\E_{\underline e\mid\underline y}G_{\underline e}(R,F),
 \qquad \bar G_{\underline y}=T_rG_{\underline y}.
 \label{eq:average-function}
\end{align}
把用 $\sigma^{\mathrm{even}}$ 所作的类似平均记为
$A_{\underline y}$ 与 $\bar A_{\underline y}=T_rA_{\underline y}$。
这些函数的绝对值以 1 为界。此外，
\[
 G_{\underline y}(R,F)=\phi_{\underline y}(W(R,F)),\qquad
 \phi_{\underline y}(w)=\E_{\underline e\mid\underline y}
 \sigma^o_{\underline x(\underline e)}(w\circ\pi_{\underline e}),
\]
其中 $\phi_{\underline y}$ 是奇的，且不依赖任何切片或行数据。于是由
\Cref{lem:tree-correlation} 之后的观察，对每个固定的
$\underline y,\mathcal D,R,F_{\mathrm{out}}$，$\bar G_{\underline y}$
在切片表项上均值为零。后文的抽取论证（\Cref{prop:extraction}）恰好
适用于这种有界奇性的生成词函数。

\begin{lemma}[稳定性恒等式]\label{lem:cut-stability}
对构造出的图的每个割，
\begin{equation}\label{eq:cut-stability}
 \cut(\sigma)=\frac12+
 \frac12\E\Stab_t^{F_{\mathcal D}}(\bar G_{\underline y})
 -\frac12\E\Stab_t^{F_{\mathcal D}}(\bar A_{\underline y}),
\end{equation}
其中两个期望都取遍 $\underline y,\mathcal D,R,F_{\mathrm{out}}$。
特别地，
\begin{equation}\label{eq:cut-upper-stability}
 \cut(\sigma)\le\frac12+
 \frac12\E\Stab_t^{F_{\mathcal D}}(\bar G_{\underline y}).
\end{equation}
\end{lemma}

\begin{proof}
以 $\underline y,\mathcal D,R,F_{\mathrm{out}}$ 为条件。采样到的边
元组与两次行重采样在两个副本之间独立。对它们取平均后，作为
$F_{\mathcal D}$ 的函数，第一个符号的期望是
$\bar A_{\underline y}+\bar G_{\underline y}$。第二个被查询词前的负
号使第二个符号在 $F'_{\mathcal D}$ 处的期望为
$\bar A_{\underline y}-\bar G_{\underline y}$。两个切片输入是普通的
$t$-相关均匀立方输入。这些函数中的行是公共潜数组 $R$；两次独立的行
重采样已被吸收进 $T_r$。把每个切片表项取反，$\bar A_{\underline y}$
不变而 $\bar G_{\underline y}$ 变号。因此二者的傅里叶支撑的总奇偶性
相反，交叉项消失。期望符号积为
$\Stab_t(\bar A_{\underline y})-\Stab_t(\bar G_{\underline y})$。
一条边以「二分之一减符号积的二分之一」的概率被割，给出
\eqref{eq:cut-stability}。$t>0$ 时稳定性非负给出
\eqref{eq:cut-upper-stability}。
\end{proof}

若某割的值至少为 $(1+B(t))/2+\delta/2$，则其奇平均的期望切片稳定性
至少为 $B(t)+\delta$。接下来三节表明：这一超额稳定性产生一个承载
可观傅里叶质量的标记，且常数与博弈维数及重复次数无关；形式化陈述是
\Cref{prop:extraction}。\Cref{sec:decoding} 再把该质量转换为对原
博弈单个副本的染色。

\section{树码的傅里叶结构}\label{sec:tree-fourier}

割依赖生成的词，因此每个保持该词的变换同样约束割的行傅里叶系数。
在每个门处，这些约束给出所有子节点对共享的一个公共傅里叶指标。
把门表取反即可区分该指标是否为求值标记。随后我们证明：在随机切片
处，把整个公共指标只保留在每对的一个子节点中，平滑后的傅里叶质量
损失很小。最后我们把保留下来的指标表示成在其他行固定后仍然自由的
行坐标。这就是 \Cref{sec:extraction} 使用的局部形式。
本节通篇假设 $f$ 是 \eqref{eq:function-of-code} 意义下的码的函数。

\subsection{行傅里叶支撑}

对每个节点 $v$，递归地定义其\emph{生成子} $P(v)$（一个叶子集）：
叶子处 $P(v)=\{v\}$，内部节点处
$P(v)=P(v_{1,0})\cup P(v_{1,1})$。
即生成子总是沿第一对子节点的两侧；与切片不同，它不含随机选择。
把同一个 $h\in H$ 加到 $P(v)$ 中的所有行上，会把该子树在 $a$ 处的
输出乘以 $(-1)^{h(a)}$。这由 Equation~\eqref{eq:gate-rule} 归纳
可得：把一对的两个子节点乘以同一符号，所有门输入 $z_g$ 不变，而
前因子乘以该符号。

对行傅里叶指标 $C=(q_i)_{i=1}^k\in(H^*)^k$，令
\[
 \chi_C(R)=\prod_i\chi_{q_i}(R_i),\qquad
 |C|=|\{i:q_i\ne0\}|,\qquad
 q_v=\sum_{i\in P(v)}q_i.
\]
把 $v$ 的子节点 $(g,b)$ 的指标记为 $q_{v,gb}$。按定义有
$q_v=q_{v,10}+q_{v,11}$。在固定表处，行展开及其平滑为
\begin{equation}
 f(R,F)=\sum_C\widehat f_C(F)\chi_C(R),\qquad
 \widehat{\bar f}_C(F)=r^{|C|}\widehat f_C(F).
 \label{eq:row-smoothing}
\end{equation}

\begin{lemma}[行对称性]
\label{lem:row-symmetry}
设 $f$ 是码的函数。$f$ 的每个非零行傅里叶系数在每个内部节点 $v$ 处
满足
\begin{equation}
 q_{v,g0}+q_{v,g1}=q_v\qquad(g\in[s]).
 \label{eq:row-relations}
\end{equation}
对表集合 $F$，记 $F^{v,-}$ 为仅把 $F_v$ 取反，$F^{v,z}$ 为把
$F_v(u)$ 换成 $F_v(u+z)$，其余表不变。则
\begin{align}
 \widehat f_C(F^{v,-})
   &=(-1)^{q_v(\one)}\widehat f_C(F),
   \label{eq:table-flip-action}\\
 \widehat f_C(F^{v,z})
   &=(-1)^{\sum_g z_g q_{v,g1}(\one)}\widehat f_C(F).
   \label{eq:table-translation-action}
\end{align}
以上陈述在行平滑后仍然成立。它们也被行傅里叶投影，以及到某个指定
整门表集合的表项总傅里叶次数范围的投影所保持。
\end{lemma}

\begin{proof}
对每个 $g$，把对 $g$ 中两个子节点的生成子都平移 $h$，对 $v$ 处子树
的效果与把 $P(v)$ 平移 $h$ 完全相同。因此两个行平移对整张词的效果
相同。取行傅里叶系数得
\[
 \chi_{q_{v,g0}+q_{v,g1}}(h)\widehat f_C(F)
   =\chi_{q_v}(h)\widehat f_C(F)\qquad(h\in H).
\]
不同特征是不同函数，证得 Equation~\eqref{eq:row-relations}。

把 $F_v$ 取反，其效果等同于把 $\one$ 加到 $P(v)$ 的所有行上。这证
明 Equation~\eqref{eq:table-flip-action}。为把 $F_v$ 的自变量平移
$z$，在每个 $z_g=1$ 的对中，把其第二个子节点的生成子加 $\one$。
这恰好改变那些门输入比特，而前因子不变。傅里叶变换给出
Equation~\eqref{eq:table-translation-action}。

行平滑与行傅里叶投影在 $C$ 上是对角的，故保持三个恒等式。在表项
符号的立方上，定义域平移置换坐标，而把门表取反则给每个单项式乘以
其奇偶符号。两种运算都保持任一指定整门表集合中的总次数，因此与
所述次数投影交换。
\end{proof}

\subsection{随机切片上的不分裂投影}

Equation~\eqref{eq:row-relations} 给每个子节点对指派同一和
$q_v$。Equation~\eqref{eq:table-flip-action} 由此有一个有用的解读：
在把 $F_v$ 取反下为奇的系数满足 $q_v(\one)=1$，故由
\eqref{eq:labels-dual} 其公共指标是一个标记。但单个子节点指标仍可以
是和为 $q_v$ 的任意泛函。

我们在一个门处化简这最后的自由度。把 $q_v$ 保留在每对的一个子节点
中，只剩下选哪个子节点的自由，而不是两个和给定的任意泛函。称 $C$
\emph{在 $v$ 处不分裂}，若对每个 $g$，有序对 $(q_{v,g0},q_{v,g1})$
为 $(q_v,0)$ 或 $(0,q_v)$。记 $S_v$ 为到这类指标的正交行傅里叶
投影。当 $q_v=0$ 时，这要求每个子节点指标为零。下一个界把代价对
切片深度取平均。它不要求指标在树的每个门处都不分裂。

\begin{lemma}[不分裂投影的平均代价]
\label{lem:unsplit}
设 $f$ 是满足 $|f|\le1$ 的码的函数，$r=1-\mu/d$（$0<\mu<1/2$），
$\bar f=T_rf$。选取与 $f$ 独立的随机切片。则
\begin{equation}
 \E_{J_0,\mathcal D}\sum_{v\in\mathcal D}
       \norm{(I-S_v)\bar f}_{2,R,F}^{2}
       \le \frac{s}{2e m\mu}.
 \label{eq:unsplit-bound}
\end{equation}
在决定 $f$ 的额外背景数据上取平均后，同样的界成立，只要切片的选取
独立于那些数据。
\end{lemma}

\begin{proof}
固定满足 Equation~\eqref{eq:row-relations} 的指标 $C$，令 $N_j$ 为
深度 $j$ 处指标 $q_v$ 非零的节点数。若 $q_v\ne0$，每个子节点对至少
含一个非零指标；该节点处不分裂条件的一次失败至少再提供一个非零
子节点。若 $q_v=0$，一次失败提供至少两个非零子节点，而其最小贡献
为零。因此深度 $j$ 处失败的个数 $B_j$ 满足
\[
 B_j\le N_{j+1}-sN_j.
\]
用切片包含概率与裂项相消得
\begin{align*}
 \E_{J_0,\mathcal D}\sum_{v\in\mathcal D}
          \one_{\{C\text{ 在 }v\text{ 处分裂}\}}
 &\le \frac1m\sum_{j=0}^{m-1}
          \frac{N_{j+1}-sN_j}{s^j}\\
 &=\frac{s}{m}\left(\frac{N_m}{s^m}-N_0\right)
 \le\frac{s|C|}{md}.
\end{align*}
此处 $N_m=|C|$。由 Equation~\eqref{eq:row-smoothing} 与 Parseval，
Equation~\eqref{eq:unsplit-bound} 的左边至多为
\[
 \frac{s}{m}\sum_C
   \left(r^{2|C|}\frac{|C|}{d}\right)
           \E_F|\widehat f_C(F)|^2.
\]
对 $x\ge0$ 有 $xe^{-2\mu x}\le(2e\mu)^{-1}$，且 $\log r\le-\mu/d$。
故每个括号因子至多 $(2e\mu)^{-1}$，而
$\sum_C\E_F|\widehat f_C(F)|^2=\norm f_2^2\le1$。
界得证。证明逐点依赖于背景数据，故其平均满足同一不等式。
\end{proof}

\subsection{自由平移与剩余行坐标}

抽取论证必须在固定大部分行之后使用公共指标。因此我们选取这样的
坐标：平移一个子节点的生成子只改变一个自由行，而其余坐标不动。
这样不分裂投影便在每个固定上下文处单独作用，且其保留的特征具有
显式形式。

固定切片 $\mathcal D$ 与节点 $v\in\mathcal D$。$2s$ 个集合
$P(v_{g,b})$ 非空且两两不交。按只依赖树的规则从每个集合中选一个
代表叶 $i_{g,b}$，令 $L_{g,b}=R_{i_{g,b}}$。对该生成子中的其他每个
叶子 $i$，把 $R_i$ 换成 $R_i+L_{g,b}$。这些生成子之外的行保持不变。
把所得的 $k-2s$ 个剩余坐标记为 $\Gamma\in H^{k-2s}$。

\begin{lemma}[自由平移坐标]
\label{lem:shift-coordinates}
坐标变换
\[
 H^k\longrightarrow H^{2s}\times H^{k-2s},\qquad
 R\longmapsto((L_{g,b}),\Gamma)
\]
是线性双射。特别地，在均匀行下，自由平移在给定 $\Gamma$ 时为独立
均匀。与完整行指标 $C$ 对应的 $L_{g,b}$ 的傅里叶指标是
$Q_{g,b}=q_{v,g b}$。

因此，在每个固定的 $\Gamma$ 处，算子 $S_v$ 在 $2s$ 个平移上的作用
恰是到不分裂平移指标的傅里叶投影。对码的函数或其行平滑，除非和
$Q_{g,0}+Q_{g,1}$ 对每个 $g$ 相同，否则平移系数为零。
Equations~\eqref{eq:table-flip-action} 与
\eqref{eq:table-translation-action} 的表作用在该固定上下文处对每个
平移系数成立，其中 $q_v=Q_{1,0}+Q_{1,1}$，$q_{v,g1}=Q_{g,1}$。
\end{lemma}

\begin{proof}
逆映射把代表行取为 $L_{g,b}$，并把其记录的差加 $L_{g,b}$ 恢复生成
子中的每一行。这证明双射性。在行特征中，$L_{g,b}$ 的系数恰是
$\sum_{i\in P(v_{g,b})}q_i=q_{v,g b}$。于是定义 $S_v$ 的条件只依赖
平移指标 $Q$，不依赖 $\Gamma$ 的任何傅里叶指标。把完整傅里叶展开
按 $Q$ 分组，即得其在每个固定纤维上的所述作用。同一分组证得支撑
与表作用恒等式：给定 $Q$ 的所有完整行系数满足同样的恒等式并带有
同样的符号。
\end{proof}

于是对非零公共指标 $q$，保留的平移指标由子集 $I\subseteq[s]$
（$q$ 出现在其中哪些对的第二个子节点）刻画。其特征为
\[
 \prod_{g\notin I}\chi_q(L_{g0})
 \prod_{g\in I}\chi_q(L_{g1}).
\]
若 $q=0$，只有全零的不分裂指标，其特征为 $1$。把 $F_v$ 取反会给
相应系数乘上 $(-1)^{q(\one)}$。在 $F_v$ 的奇部上，$q$ 是标记因而
非零，故由 $q$ 与 $I$ 的刻画唯一。对这样的标记，把 $F_v$ 的自变量
平移 $z\in\F^s$ 会给其系数乘上 $(-1)^{\sum_{g\in I}z_g}$。
因此同一子集 $I$ 同时记录行特征与表平移特征。这就是
\Cref{sec:extraction} 使用的局部傅里叶描述，在每个固定上下文处
有效。

\section{向量值哈希引理}\label{sec:hashing}

下面的引理把特征之间的弥散转化为表项之间的弥散。其常数与特征空间
的维数以及被哈希向量的维数无关。在应用中，$q$ 将索引候选标记，
$z$ 将索引门表的表项。所得的哈希向量将收集对一个表项的影响有贡献
的全部傅里叶系数，故其范数平方即为该影响。这就是我们需要向量值
命题而非标量估计的原因。

\begin{lemma}[向量值哈希]\label{lem:hashing}
设 $H$ 是 $\F$ 上的有限维向量空间，对 $q\in H^*$ 与 $u\in H$ 记
$\chi_q(u)=(-1)^{q(u)}$；设 $\mathcal V$ 是实希尔伯特空间，
$s\ge1$ 为整数，且
\[
 c(q,z)\in\mathcal V,
 \qquad q\in H^*,\quad z\in\F^s.
\]
令
\[
 w_q=\sum_{z\in\F^s}\norm{c(q,z)}^2,
 \qquad M=\sum_{q\in H^*}w_q.
\]
独立均匀地选取 $L,J_1,\ldots,J_s\in H$，并定义
\[
 \mathbf J(q)=(q(J_1),\ldots,q(J_s)),
 \qquad
 U(z)=\sum_{q\in H^*}\chi_q(L)c(q,z+\mathbf J(q)).
\]
则
\begin{equation}\label{eq:hashing-total-mass}
 \E\sum_{z\in\F^s}\norm{U(z)}^2=M.
\end{equation}
对每个 $\zeta>0$，存在只依赖 $\zeta$ 的 $\gamma>0$ 与
$s_0\in\mathbb N$，使得对每个 $s\ge s_0$，
\begin{equation}\label{eq:hashing-dispersion}
 \max_{q\in H^*}w_q\le\gamma M
 \quad\Longrightarrow\quad
 \E\max_{z\in\F^s}\norm{U(z)}^2\le\zeta M.
\end{equation}
\end{lemma}

\begin{proof}
下文对 $H$ 与 $\F^s$ 的期望都使用均匀概率测度；定义 $w_q$ 与 $M$
的显示求和不作归一化。在给定 $J_g$ 的条件下，$L$ 的特征的正交性给出
\[
 \E_L\norm{U(z)}^2
 =\sum_q\norm{c(q,z+\mathbf J(q))}^2.
\]
对 $z$ 求和即证 \eqref{eq:hashing-total-mass}。若 $M=0$，则每个向量
$c(q,z)$ 为零，\eqref{eq:hashing-dispersion} 平凡成立。以下设
$M>0$。

令 $h=2^s$，$\rho=\max_q w_q/M$。下面对每个 $C,D>0$ 证明定量界，
再选取其参数。主要估计把最大范数平方用四次幂之和的平方根来控制。
经傅里叶换元后，该四阶矩中的大多数项迫使所有四个因子取同一特征
$q$；其总和由特征质量的弥散控制。两次截断将控制剩余项以及换元给
阵列带来的代价：
\begin{equation}\label{eq:hashing-quantitative}
 \E\max_z\norm{U(z)}^2
 \le 2M\left[
 \sqrt{2\rho+\frac{2C^2}{D^2}+\frac{3D^4}{h}}
 +\left(\frac1{\sqrt C}+\frac1D\right)^2+\frac1h
 \right].
\end{equation}


\paragraph{傅里叶表示。}
对 $p\in\F^s$ 与 $u\in H$，定义
\[
 \alpha_q(p)=\sum_{z\in\F^s}(-1)^{p\cdot z}c(q,z),
 \qquad
 A_p(u)=\sum_{q\in H^*}\chi_q(u)\alpha_q(p),
 \qquad
 S_p=L+\sum_{g=1}^s p_gJ_g.
\]
在 $z$ 上作傅里叶反演得
\begin{equation}\label{eq:hashing-inversion}
 U(z)=\E_p(-1)^{p\cdot z}A_p(S_p).
\end{equation}
记 $m_p=\E_u\norm{A_p(u)}^2$，两个群上的 Parseval 恒等式给出
\begin{equation}\label{eq:hashing-parseval}
 m_p=\sum_q\norm{\alpha_q(p)}^2,
 \qquad \E_p m_p=M,
 \qquad \E_p\norm{\alpha_q(p)}^2=w_q.
\end{equation}
对不同的 $p,p'$，变量 $S_p,S_{p'}$ 独立且均匀于 $H$：系数向量
$(1,p),(1,p')$ 在 $\F$ 上线性无关。同理，任意三个不同的
$p_1,p_2,p_3$ 给出三个独立均匀变量 $S_{p_1},S_{p_2},S_{p_3}$。
事实上，其系数向量的任一零和非空子集必有偶数个元素，因而由两个
相等向量组成。

\paragraph{两次截断。}
我们先丢弃平均范数平方 $m_p$ 异常大的指标 $p$，再在其余指标上丢弃
$A_p(u)$ 的异常大值。第一个截断使第二个截断下每个特征系数的改变
一致地小；第二个截断给出四阶矩指标重合时所需的逐点界。我们将用一
阶矩界与由样本 $S_p$ 两两独立得到的方差界来控制被丢弃部分。定义
\[
 A'_p(u)=
 \begin{cases}
 A_p(u),&m_p\le CM\text{ 且 }\norm{A_p(u)}\le D\sqrt M,\\
 0,&\text{否则},
 \end{cases}
\]
并令
\[
 \alpha'_q(p)=\E_u\chi_q(u)A'_p(u),
 \qquad w'_q=\E_p\norm{\alpha'_q(p)}^2,
 \qquad U'(z)=\E_p(-1)^{p\cdot z}A'_p(S_p).
\]
逐点截断不增加总范数平方，故 Parseval 给出 $\sum_qw'_q\le M$。若
$m_p\le CM$，则
\[
 \norm{\alpha_q(p)-\alpha'_q(p)}
 \le\E_u\norm{A_p(u)}
       \one_{\{\norm{A_p(u)}>D\sqrt M\}}
 \le\frac{m_p}{D\sqrt M}
 \le\frac{C\sqrt M}{D}.
\]
若 $m_p>CM$，则 $\alpha'_q(p)=0$。因此
\begin{equation}\label{eq:hashing-truncated-mass}
 \sum_qw'_q\le M,
 \qquad
 w'_q\le 2w_q+\frac{2C^2M}{D^2}.
\end{equation}

令
\[
 X_p=\norm{(A_p-A'_p)(S_p)},\qquad X=\E_pX_p.
\]
Equation~\eqref{eq:hashing-inversion} 与三角不等式蕴含
$\max_z\norm{U(z)-U'(z)}\le X$。$m_p>CM$ 的指标对 $\E X$ 的贡献
至多为
\[
 \E_p\one_{\{m_p>CM\}}\sqrt{m_p}
 \le\frac{\E_p m_p}{\sqrt{CM}}=\sqrt{M/C}.
\]
其余指标的贡献至多 $\E_p m_p/(D\sqrt M)=\sqrt M/D$。此外，由
$S_p$ 两两独立知 $X_p$ 两两独立。由于 $A_p-A'_p$ 在每点要么为
$A_p$ 要么为零，
\[
 \operatorname{Var}(X)
 =h^{-2}\sum_p\operatorname{Var}(X_p)
 \le h^{-2}\sum_p\E_u\norm{A_p(u)-A'_p(u)}^2
 \le\frac Mh.
\]
我们已证明
\begin{equation}\label{eq:hashing-truncation-error}
 \E\max_z\norm{U(z)-U'(z)}^2
 \le M\left[\left(\frac1{\sqrt C}+\frac1D\right)^2+\frac1h\right].
\end{equation}

\paragraph{四阶矩。}
为简洁记 $B_p=A'_p(S_p)$。展开范数的四次幂并对 $z$ 中的特征求和，
得精确恒等式
\begin{align}
 \E\sum_z\norm{U'(z)}^4
 &=h^{-3}\sum_{p_1,p_2,p_3}
   \E\bigl[\ip{B_{p_1}}{B_{p_2}}
                  \ip{B_{p_3}}{B_{p_1+p_2+p_3}}\bigr]\notag\\
 &=\E_{p_1,p_2,p_3}\E
   \bigl[\ip{B_{p_1}}{B_{p_2}}\ip{B_{p_3}}{B_{p_4}}\bigr],
 \qquad p_4=p_1+p_2+p_3.
 \label{eq:hashing-fourth-expansion}
\end{align}
第二行中 $p_1,p_2,p_3$ 是独立均匀指标。第一行的因子 $h^{-3}$ 来自
傅里叶反演中的四个因子 $h^{-1}$ 以及 $h$ 个特征对 $z$ 的求和。

若 $p_1,p_2,p_3,p_4$ 中任意两个重合，则 $p_1,p_2,p_3$ 中某两个
重合。该事件的概率至多 $3/h$。由 $\norm{B_p}\le D\sqrt M$，其对
\eqref{eq:hashing-fourth-expansion} 的贡献绝对值至多 $3D^4M^2/h$。

现在固定四个不同指标并设 $p_4=p_1+p_2+p_3$。前三个 $S_{p_i}$ 是
$H$ 中独立均匀元素，而 $S_{p_4}=S_{p_1}+S_{p_2}+S_{p_3}$。把四个
$A'_{p_i}$ 按特征展开并对这三个独立元素取平均，得
\begin{equation}\label{eq:hashing-common-character}
 \E\bigl[\ip{B_{p_1}}{B_{p_2}}\ip{B_{p_3}}{B_{p_4}}\bigr]
 =\sum_q
   \ip{\alpha'_q(p_1)}{\alpha'_q(p_2)}
   \ip{\alpha'_q(p_3)}{\alpha'_q(p_4)}.
\end{equation}
具体地，由 $q_1,q_2,q_3,q_4$ 索引的特征产生因子
\[
 \prod_{i=1}^3\chi_{q_i+q_4}(S_{p_i}),
\]
除非 $q_1=q_2=q_3=q_4$，其期望为零。

\eqref{eq:hashing-common-character} 的右端不必非负。为得上界，把
每个被加项换成非负的 $\prod_{i=1}^4\norm{\alpha'_q(p_i)}$，再去掉
四指标互异这一限制。对每个固定的 $q$，Cauchy--Schwarz 给出
\begin{align*}
 \E_{p_1,p_2,p_3}\prod_{i=1}^4\norm{\alpha'_q(p_i)}
 &\le
 \left(\E\norm{\alpha'_q(p_1)}^2
             \norm{\alpha'_q(p_2)}^2\right)^{1/2}
 \left(\E\norm{\alpha'_q(p_3)}^2
             \norm{\alpha'_q(p_4)}^2\right)^{1/2}\\
 &=(w'_q)^2.
\end{align*}
最后一个等式用了每对 $(p_1,p_2)$ 与 $(p_3,p_4)$ 的无条件独立均匀
律；它不断言在以互异为条件之后的独立性。于是非退化四阶矩被
$\sum_q(w'_q)^2$ 控制：特征之间的弥散就此控制表项之间的集中。把这些
估计与 \eqref{eq:hashing-truncated-mass} 结合得
\begin{align*}
 \E\sum_z\norm{U'(z)}^4
 &\le \sum_q(w'_q)^2+\frac{3D^4M^2}{h}\\
 &\le M^2\left(2\rho+\frac{2C^2}{D^2}+\frac{3D^4}{h}\right).
\end{align*}
再用一次 Cauchy--Schwarz，
\begin{equation}\label{eq:hashing-truncated-maximum}
 \E\max_z\norm{U'(z)}^2
 \le\left(\E\sum_z\norm{U'(z)}^4\right)^{1/2}
 \le M\sqrt{2\rho+\frac{2C^2}{D^2}+\frac{3D^4}{h}}.
\end{equation}
不等式 $(a+b)^2\le2a^2+2b^2$ 连同
\eqref{eq:hashing-truncation-error} 与
\eqref{eq:hashing-truncated-maximum} 证明
\eqref{eq:hashing-quantitative}。

最后令 $\varepsilon=\min\{\zeta,1\}$，并选
\[
 C=\frac{64}{\varepsilon},\qquad
 D=\frac{8C}{\varepsilon},\qquad
 \gamma=\frac{\varepsilon^2}{128},\qquad
 2^{s_0}\ge
 \max\left\{\frac{192D^4}{\varepsilon^2},\frac{16}{\varepsilon}\right\}.
\]
若 $s\ge s_0$ 且 $\rho\le\gamma$，则
\eqref{eq:hashing-quantitative} 中的平方根项至多 $\varepsilon/4$。
又 $D\ge\sqrt C$，故其括号中另两项各至多 $\varepsilon/16$。所得界
为 $3\varepsilon M/4\le\zeta M$，以只依赖 $\zeta$ 的常数建立
\eqref{eq:hashing-dispersion}。
\end{proof}

\section{从超额稳定性到标记}\label{sec:extraction}

本节给出归约所用的解析命题。其常数不依赖仿射标记空间的维数。函数
可以依赖额外的随机参数 $\omega$，其律与所有行、门表及随机切片独立。
在应用中，$\omega$ 将是博弈的右端点元组。

设 $\phi_\omega:\bits^{\mathcal A}\to[-1,1]$ 是奇的，即
$\phi_\omega(-w)=-\phi_\omega(w)$，并令
\[
 f_\omega(R,F)=\phi_\omega(W(R,F)),\qquad
 \bar f_\omega=T_r f_\omega,\qquad r=1-\mu/d.
\]
词、树与随机切片 $\mathcal D$ 均为 \Cref{sec:tree} 中者。把
$\mathcal D$ 之外的表记为 $F_{\mathrm{out}}$。一旦
$\omega,\mathcal D,R,F_{\mathrm{out}}$ 固定，$\bar f_\omega$ 便是
独立切片表比特的函数。

对节点 $v\in\mathcal D$，使用 \Cref{lem:shift-coordinates} 的自由
平移坐标 $(L_{gb})_{g\in[s],b\in\{0,1\}}$ 与剩余上下文
$\Gamma\in H^{k-2s}$。带下标 ``shifts'' 的帽子表示只对这 $2s$ 个
自由行作展开。对标记 $q\in\mathcal A$（等价地，$q\in H^*$ 且
$q(\one)=1$），定义
\begin{equation}\label{eq:mass}
 m_q(f;\Gamma,F_{\mathrm{out}})
 =\E_{F_{\mathcal D}}
 \sum_{\substack{Q\in(H^*)^{2s}\\
                  Q_{g0}+Q_{g1}=q\ (g\in[s])}}
 \abs{\widehat f^{\,\mathrm{shifts}}(Q)}^2.
\end{equation}
对 $v,\mathcal D,\omega$ 的依赖从略。此定义中的指标集随 $q$ 变化
两两不交。对固定的 $q$，该质量包含把公共和分派到每对两个子节点的
一切方式。我们要为原平滑函数 $\bar f_\omega$ 寻找这样的质量。
不分裂投影与表次数截断只用于找出承载该质量的标记；两个投影都不是
结论的一部分。保留具有给定公共和的全部指标，将使我们在
\Cref{sec:decoding} 中能在仿射拉回下比较质量。

特别地，若 $|f|\le1$，则在每个固定上下文处
\begin{equation}\label{eq:mass-bound}
 0\le\sum_{q\in\mathcal A}m_q(f;\Gamma,F_{\mathrm{out}})\le1.
\end{equation}

\begin{proposition}[维数无关的抽取]\label{prop:extraction}
固定 $t\in(0,1)$、$0<\delta<1-B(t)$ 与 $0<\mu<1/2$。
存在只依赖 $t,\delta,\mu$ 的树参数 $s\ge2,m\ge1$ 与常数
$p_*,\theta>0$，具有以下性质。对任意仿射标记空间 $\mathcal A$ 与
如上任意一族 $\phi_\omega$，若
\begin{equation}\label{eq:excess-stability}
 \E_{\omega,\mathcal D,R,F_{\mathrm{out}}}
 \Stab_t^{F_{\mathcal D}}(\bar f_\omega)
 \ge B(t)+\delta,
\end{equation}
则
\begin{equation}\label{eq:extraction-event}
 \Pr\left[
   \max_{q\in\mathcal A}
   m_q(\bar f_\omega;\Gamma,F_{\mathrm{out}})\ge\theta
 \right]\ge p_*.
\end{equation}
这里采样 $\omega$、随机切片 $\mathcal D$ 与均匀节点 $v\in\mathcal D$；
在其条件下，独立均匀地采样 $F_{\mathrm{out}}$ 与 $\Gamma$。
\end{proposition}

\begin{proof}
超额稳定性首先产生一个有影响的切片表项。我们在不分裂投影之后保留
该影响，丢弃其门表中偶性的部分，并把剩余影响表示为由标记索引的
向量的仿射哈希。哈希引理随后迫使某个标记承载可观的条件下质量。

下文选取的所有常数都与 $\mathcal A$ 及 $\omega$ 的律无关。先选
次数截断 $K$ 与影响阈值 $\tau$。估计过程中暂不指定 $s,m$：
第 4 步将选 $s$ 以控制偶性与哈希贡献，再选 $m$ 以控制不分裂投影的
代价。

\paragraph{第 1 步：找出有影响的切片表项。}
把切片表同时取反会把词取反。由于 $\phi_\omega$ 是奇的且平滑保持该
性质，$\bar f_\omega$ 在每个固定的
$\omega,\mathcal D,R,F_{\mathrm{out}}$ 处于切片比特上均值为零。
它也落在 $[-1,1]$ 中。以 $\xi=\delta/2$ 应用 \Cref{lem:mis}，得
$K\ge1$ 与 $\tau>0$。不妨缩小 $\tau$ 以保证 $\tau\le1$。
由于稳定性至多为 1，\eqref{eq:excess-stability} 蕴含：以至少
\[
 \eta=\delta/2
\]
的概率，某个 $v\in\mathcal D$ 的某表项 $z\in\F^s$ 满足
$\Inf_{v,z}^{\le K}(\bar f_\omega)\ge\tau$。
事实上，若该事件的概率为 $p$，则期望稳定性至多
$B(t)+\delta/2+p$。

设 $S_v$ 为到在 $v$ 处不分裂的指标的正交行投影。令
\[
 H_v=P_{\le K}^{F_{\mathcal D}}S_v\bar f_\omega.
\]
由 \Cref{lem:unsplit}，投影损失所需的界是
\begin{equation}\label{eq:choose-m}
 \E\sum_{v\in\mathcal D}
 \norm{(1-S_v)\bar f_\omega}_{2,R,F}^2
 \le\frac{s}{2em\mu}\le\frac{\eta\tau}{8}.
\end{equation}
对任意固定的 $s$，增大 $m$ 即可保证最后一个不等式。同一期望也是
相应的 $F_{\mathcal D}$ 条件范数平方之和在
$\omega,\mathcal D,R,F_{\mathrm{out}}$ 上的期望。Markov 不等式
表明：除一个概率 $\eta/2$ 的事件外，该和至多 $\tau/4$。在与有影响
表项事件的交上，定义该影响的投影的三角不等式给出
\[
 \sqrt{\Inf_{v,z}(H_v)}
 \ge\sqrt{\Inf_{v,z}^{\le K}(\bar f_\omega)}
      -\norm{(1-S_v)\bar f_\omega}_{2,F_{\mathcal D}}
 \ge\frac{\sqrt\tau}{2}.
\]
因此
\begin{equation}\label{eq:large-max-sum}
 \E\sum_{v\in\mathcal D}\max_z\Inf_{v,z}(H_v)
 \ge\frac{\eta\tau}{8}=:\Lambda.
\end{equation}

我们还要用
\begin{equation}\label{eq:total-influence}
 \E\sum_{v\in\mathcal D}\sum_z\Inf_{v,z}(H_v)\le K.
\end{equation}
验证如下：对全部行积分。每个 $S_v$ 是行中的正交投影，与切片比特的
傅里叶投影交换，故只会减少其表项处的积分影响。逐个去掉 $S_v$ 后，
左边以 $P_{\le K}^{F_{\mathcal D}}\bar f_\omega$ 的总影响对其余
变量的平均为界。Equation~\eqref{eq:degree-influence} 与
$|\bar f_\omega|\le1$ 给出 \eqref{eq:total-influence}。
不断言也不需要 $H_v$ 的逐点有界性。

\paragraph{第 2 步：丢弃对某张门的偶性依赖。}
固定 $v\in\mathcal D$，把 $H_v$ 在仅把 $F_v$ 取反下分解为偶部与
奇部。两部分在切片比特上的傅里叶支撑不交，在保留含任一指定表项的
单项式后亦然。于是
\[
 \Inf_{v,z}(H_v)=
 \Inf_{v,z}(H_v^{\mathrm{even}})
 +\Inf_{v,z}(H_v^{\mathrm{odd}}).
\]
由 \Cref{lem:row-symmetry} 的符号作用，偶部在行支撑上满足
$q_v(\one)=0$。由于指标不分裂，其所有子节点指标要么为 $0$ 要么为
$q_v$，故在 $\one$ 上也取零。该引理中的定义域平移作用表明：在行与
切片外的表固定时，偶部在 $F_v$ 的定义域平移下不变。这些平移在其
表项上传递。因此 $v$ 处的所有表项影响相等。由
\eqref{eq:total-influence}，
\begin{equation}\label{eq:even-influence}
 \E\sum_{v\in\mathcal D}
 \max_z\Inf_{v,z}(H_v^{\mathrm{even}})\le\frac{K}{2^s}.
\end{equation}

\paragraph{第 3 步：把奇影响表示为仿射哈希。}
固定 $\omega,\mathcal D,v,\Gamma,F_{\mathrm{out}}$。剩余行坐标
$L_{gb}$ 是 $H$ 中独立均匀元素。在奇部上，公共对和 $q$ 由
\eqref{eq:labels-dual} 是标记，因而非零。具有公共和 $q$ 的不分裂
平移指标在子集 $I\subseteq[s]$ 上把 $q$ 放在第二个子节点、其余放在
第一个子节点。记其在 $H_v^{\mathrm{odd}}$ 中的系数为
$a_{q,I}(F_{\mathcal D})$。

令
\[
 L_\Sigma=\sum_{g=1}^sL_{g0},\qquad J_g=L_{g0}+L_{g1},\qquad
 \mathbf J(q)=(q(J_1),\ldots,q(J_s)).
\]
$s+1$ 个型 $L_\Sigma,J_1,\ldots,J_s$ 独立且均匀：所示的从 $H^{2s}$
到 $H^{s+1}$ 的线性映射是满射。其余 $s-1$ 个坐标不出现在不分裂
展开中。平移指标 $(q,I)$ 的特征是
\[
 \chi_q(L_\Sigma)(-1)^{\sum_{g\in I}q(J_g)}.
\]
另一方面，定义域平移恒等式给出
\[
 a_{q,I}(F_v(\cdot+z),F_{\mathcal D\setminus\{v\}})
 =(-1)^{\sum_{g\in I}z_g}a_{q,I}(F_{\mathcal D}).
\]
令 $\Psi_q=\sum_{I\subseteq[s]}a_{q,I}$，我们得到
\begin{equation}\label{eq:odd-hash-form}
 H_v^{\mathrm{odd}}
 =\sum_{q\in\mathcal A}\chi_q(L_\Sigma)
 \Psi_q\bigl(F_v(\cdot+\mathbf J(q)),
                   F_{\mathcal D\setminus\{v\}}\bigr).
\end{equation}
不同 $I$ 的函数 $a_{q,I}$ 在均匀切片比特上正交：两个特征不同的定义
域平移保持其内积并改变其符号。因此
\begin{equation}\label{eq:psi-mass}
 \norm{\Psi_q}_{2,F_{\mathcal D}}^2
 =\sum_I\norm{a_{q,I}}_{2,F_{\mathcal D}}^2
 \le m_q(\bar f_\omega;\Gamma,F_{\mathrm{out}}).
\end{equation}
该不等式成立是因为不分裂行投影由自由平移指标决定，而次数与奇偶
投影在每个固定上下文处是切片比特中的正交投影。每个 $\Psi_q$ 的
次数至多为 $K$。

Equation~\eqref{eq:odd-hash-form} 已把自由平移特征转换为门表的
平移。接下来我们把它的布尔傅里叶系数收集成范数平方即表项影响的
向量。这些傅里叶系数由表项集合索引；此前的指标 $I\subseteq[s]$
刻画的则是表的定义域 $\F^s$ 的一个特征。

把门 $v$ 处表项集 $B\subseteq\F^s$、其余切片门表项集 $D$ 的切片
傅里叶系数记为 $\widehat\Psi_q(B,D)$。对 $z\in\F^s$ 定义实向量
$c(q,z)$，其坐标由对 $(A,D)$（$0\in A\subseteq\F^s$）索引：
\[
 c(q,z)_{A,D}=\widehat\Psi_q(z+A,D).
\]
用相对于 $z$ 的表项集，使该向量空间对每个 $z$ 与每个标记 $q$ 都
相同。对非标记的泛函置 $c(q,z)=0$。对 $z$ 处的影响，含 $z$ 的切片
傅里叶集 $B$ 有唯一形式 $B=z+A$（$0\in A$）。平移
$\mathbf J(q)$ 把其系数送到 $\widehat\Psi_q(z+A+\mathbf J(q),D)$。
于是恰好有
\begin{equation}\label{eq:influence-hash}
 \Inf_{v,z}(H_v^{\mathrm{odd}})
 =\norm*{\sum_q\chi_q(L_\Sigma)c(q,z+\mathbf J(q))}^2.
\end{equation}
记 $w_q=\sum_z\norm{c(q,z)}^2$ 与 $M_v=\sum_qw_q$，则
\begin{equation}\label{eq:wq-mq}
 w_q=\sum_{B,D}|B|\abs{\widehat\Psi_q(B,D)}^2
 \le K\norm{\Psi_q}_2^2
 \le K m_q(\bar f_\omega;\Gamma,F_{\mathrm{out}}).
\end{equation}
因子 $|B|$ 计其影响含该傅里叶单项式的表项个数；它至多单项式的总
次数，从而至多 $K$。
于是每个上下文处 $M_v\le K$。\Cref{lem:hashing} 的总质量恒等式与
\eqref{eq:influence-hash} 一并给出
\begin{equation}\label{eq:sum-M}
 \E\sum_{v\in\mathcal D}M_v
 =\E\sum_{v\in\mathcal D}\sum_z
     \Inf_{v,z}(H_v^{\mathrm{odd}})\le K.
\end{equation}
对该恒等式，每个节点的自由坐标分开积分；它们不同的坐标系不影响
均匀行律。

\paragraph{第 4 步：迫使出现重标记。}
选 $\zeta>0$ 使 $\zeta K\le\Lambda/4$。应用 \Cref{lem:hashing}，
得 $\gamma>0$ 与 $s$ 的下界。选足够大的 $s\ge2$ 满足该引理与
$K2^{-s}\le\Lambda/4$。再按 \eqref{eq:choose-m} 选 $m$。最后令
\begin{equation}\label{eq:extraction-constants}
 M_0=\frac{\Lambda}{4\cdot2^m},\qquad
 \theta=\frac{\gamma M_0}{K},\qquad
 p_* =\frac{\Lambda}{4K\,2^m}.
\end{equation}

对每个节点，以 $\omega,\mathcal D,v,\Gamma,F_{\mathrm{out}}$ 为
条件并对自由平移积分。把这些上下文分成三类：
\begin{enumerate}[label=(\roman*)]
\item 对每个 $q$ 有 $w_q\le\gamma M_v$；
\item 第一条不成立且 $M_v<M_0$；
\item $M_v\ge M_0$ 且存在 $q$ 使 $w_q>\gamma M_v$。
\end{enumerate}
在类 (i) 上，哈希引理把期望最大奇影响界定为 $\zeta M_v$。求和并用
\eqref{eq:sum-M}，代价至多 $\Lambda/4$。在类 (ii) 上，期望意义的
最大值以总影响即 $M_v<M_0$ 为界。由于 $|\mathcal D|\le2^m$，这些
上下文代价至多 $\Lambda/4$。在类 (iii) 上，同一总质量界对每节点
至多为 $K$。
Equation~\eqref{eq:even-influence} 把所有偶性贡献界定在 $\Lambda/4$。
设 $N_{\mathrm{iii}}$ 计其导出上下文落在类 (iii) 的节点个数。把
四项贡献与 \eqref{eq:large-max-sum} 比较：
\[
 \Lambda
 \le K2^{-s}+\zeta K+2^m M_0+K\E N_{\mathrm{iii}}
 \le\frac{3\Lambda}{4}+K\E N_{\mathrm{iii}}.
\]
于是 $\E N_{\mathrm{iii}}\ge\Lambda/(4K)$。在每个这样的节点处，
\eqref{eq:wq-mq} 给出标记使 $m_q>\gamma M_0/K=\theta$。
从切片中均匀选节点至多损失因子 $2^m$，得
\eqref{eq:extraction-event}。命题得证。
\end{proof}

\section{带隐藏坐标的解码}
\label{sec:decoding}

现在把 Proposition~\ref{prop:extraction} 提供的傅里叶质量转换为对
原仿射博弈的染色。困难在于：质量是在固定行上下文之后度量的。
$X$ 顶点处的染色者不知道其关联边上的投影，因而通常无法把该上下文
提升到自己的标记空间。我们的解决办法是：重复博弈坐标，并植入一个
坐标，使上下文与潜在重采样行在其上的线性系数为零。把被植入坐标
隐藏起来，就使所得分布接近原分布。跨越未知投影的比较于是只涉及
$2s$ 个自由平移，其损失与 $n$ 和 $T$ 无关。

本节通篇固定割表。我们沿用前几节的树参数与行坐标。对切片
$\mathcal D$ 与节点 $v\in\mathcal D$ 的每组选取，按确定性的顺序选取
代表叶，固定坐标变换
\[
 R\longleftrightarrow (L,\Gamma),\qquad
 L=(L_{gb})_{g\in[s],\,b\in\{0,1\}}\in H^{2s},
 \quad \Gamma\in H^{k-2s}.
\]
该变换只用行坐标的加法，且只依赖树与 $v$。特别地，它与仿射型的
拉回交换。质量 $m_q$ 总是指这些坐标以及 Equation~\eqref{eq:mass}
中对切片表的均匀平均。

\subsection{去掉行平滑}

设 $R^{\mathrm{new}}_1,\ldots,R^{\mathrm{new}}_k$ 是 $H$ 中独立均匀
元素，并独立地设 $B_1,\ldots,B_k$ 为 $\Pr(B_i=1)=1-r$ 的 Bernoulli
变量。令
\[
 E_i=B_iR^{\mathrm{new}}_i,
 \qquad E=(E_1,\ldots,E_k).
\]
对任意固定的行数组 $R$，数组 $R+E$ 具有测试所用的单副本行重采样
分布。把 $E$ 在上述线性变换下的坐标记为 $(E_L,E_\Gamma)$。

\begin{lemma}[去掉平滑]\label{lem:unsmooth}
对每组固定的 $\underline y$、$\mathcal D$、$v$、$\Gamma$、
$F_{\mathrm{out}}$ 与标记 $q\in\mathcal A$，
\begin{equation}\label{eq:unsmooth}
 m_q(\bar G_{\underline y};\Gamma,F_{\mathrm{out}})
 \leq
 \E_{E,\,\underline e\mid\underline y}
 m_q(G_{\underline e};\Gamma+E_\Gamma,F_{\mathrm{out}}).
\end{equation}
此处噪声 $E$ 与条件采样的边元组 $\underline e$ 独立。
\end{lemma}

\begin{proof}
对该固定上下文，把 $(L,F_{\mathcal D})$ 的函数视为带均匀概率测度的
实希尔伯特空间中的元素。令 $\Pi_q$ 为到满足
$Q_{g0}+Q_{g1}=q$（对所有 $g$）的平移傅里叶指标 $Q$ 的正交投影。
则相应的质量即该投影的范数平方。恒等式
\[
 \bar G_{\underline y}(L,\Gamma,F)
 =\E_{E,\,\underline e\mid\underline y}
 G_{\underline e}(L+E_L,\Gamma+E_\Gamma,F)
\]
连同希尔伯特空间范数平方的凸性给出所需不等式。事实上，平移
$E_L$ 与 $\Pi_q$ 交换且保持其范数：在每个平移傅里叶系数上它只乘以
一个符号。
\end{proof}

一劳永逸地为每个标记空间选定一个排序，并用该排序打破平局，定义
\begin{equation}\label{eq:decoder-choice}
 q^*=q^*(\underline y,\mathcal D,v,\Gamma,F_{\mathrm{out}})
 \in\operatorname*{arg\,max}_{q\in\mathcal A}
 m_q(\bar G_{\underline y};\Gamma,F_{\mathrm{out}}).
\end{equation}
该定义使用原始的条件边平均与原始行平滑算子。在下文修改行数据的
分布时它保持不变。

按 Proposition~\ref{prop:extraction} 采样 $\underline y$、切片、其
均匀节点、上下文与切片外的表，然后采样 $R^{\mathrm{new}}$、$B$ 与
$\underline e\mid\underline y$。把该命题与 Lemma~\ref{lem:unsmooth}
结合给出
\begin{equation}\label{eq:unplanted-score}
 \E\,m_{q^*}(G_{\underline e};
              \Gamma+E_\Gamma,F_{\mathrm{out}})
 \geq p_*\theta.
\end{equation}
此期望中的质量落在 $[0,1]$ 内，因为 $G_{\underline e}$ 的绝对值以
1 为界。

\subsection{植入一个坐标}

把每条仿射行写成 $\mathcal A=(\F^n)^T$ 上的坐标。除单个常数系数
外，它对 $T$ 个槽位各有一块 $n$ 个线性系数。对槽位 $a\in[T]$，令
$Z_a$ 为槽位 $a$ 的块在所有 $k-2s$ 个上下文行与所有 $k$ 个潜在行
$R^{\mathrm{new}}_i$ 中为零的事件。对常数系数不施加任何条件。在
原始的、未植入的律下，
\begin{equation}\label{eq:zero-slot-probability}
 p_0:=\Pr(Z_a)=2^{-n(2k-2s)}.
\end{equation}
在给定深度、切片与节点的条件下，事件 $Z_a$ 相互独立：它们涉及
独立均匀系数的不相交块。

\emph{植入实验}选取均匀标记 $a_*\in[T]$，并把 $a_*$ 处的这些块置
零。其余行系数、噪声掩码、切片外的表与博弈边都保持原分布。等价地，
在 $a_*=a$ 的条件下，这是以 $Z_a$ 为条件的未植入律。

\begin{lemma}[隐藏被植入的坐标]\label{lem:planting}
设 $\mathsf U$ 为未植入律，$\mathsf P$ 为遗忘标记后的植入律。两个
律都取在下述完整数据上：
\[
 \bigl(\underline e,J_0,\mathcal D,v,F_{\mathrm{out}},
        \Gamma,R^{\mathrm{new}},B\bigr).
\]
则
\begin{equation}\label{eq:planting-tv}
 \TV(\mathsf P,\mathsf U)
 \leq \frac12\sqrt{\frac{1-p_0}{Tp_0}}
 \leq \frac{1}{2\sqrt{Tp_0}}.
\end{equation}
因此，若
\begin{equation}\label{eq:decoder-repetition}
 T\geq\frac{1}{p_0(p_*\theta)^2},
\end{equation}
则在植入律下
\begin{equation}\label{eq:planted-score}
 \E_{\mathsf P}\,m_{q^*}(G_{\underline e};
               \Gamma+E_\Gamma,F_{\mathrm{out}})
 \geq \frac{p_*\theta}{2}.
\end{equation}
\end{lemma}

\begin{proof}
令 $Z=\sum_{a=1}^T\one_{Z_a}$。把条件密度对均匀标记取平均给出
\[
 \frac{d\mathsf P}{d\mathsf U}=\frac{Z}{Tp_0}.
\]
在给定树数据的条件下，$Z$ 具有参数 $T,p_0$ 的二项分布；这两个参数
不依赖树数据。于是 Cauchy--Schwarz 给出
\[
 \TV(\mathsf P,\mathsf U)
 =\frac12\E_{\mathsf U}\abs{\frac{Z}{Tp_0}-1}
 \leq\frac12\sqrt{\frac{1-p_0}{Tp_0}}.
\]
Equation~\eqref{eq:unplanted-score} 中的分数是所列数据（不含标记）
的 $[0,1]$ 值函数。特别地，$q^*$ 是 Equation~\eqref{eq:decoder-choice}
定义的固定函数；它不使用标记、不使用采样关联边（超出其 $Y$
端点）、也不使用采样噪声。它与零槽位的统计依赖不造成困难。该分数
在两律下的期望之差至多其全变差距离。
Equations~\eqref{eq:unplanted-score} 与 \eqref{eq:decoder-repetition}
因此蕴含 Equation~\eqref{eq:planted-score}。
\end{proof}

\subsection{仿射限制下的傅里叶质量}

下一个初等观察解释解码的损失。设 $\rho:\widetilde{\mathcal A}\to\mathcal A$
为纤维为二元的满射仿射映射，并记
\[
 \widetilde H=\Aff(\widetilde{\mathcal A},\F),\qquad
 \iota=\rho^*:H\hookrightarrow\widetilde H.
\]
即使在 $\rho$ 有非零平移部分时，拉回在仿射型向量空间上仍是线性的。
$\iota$ 的像余维数为 1，且 $\iota(\one)=\one$。其对偶限制映射
$\iota^*:\widetilde H^*\to H^*$ 纤维为二元，且在求值泛函上与
$\rho$ 一致。

\begin{lemma}[平移质量的限制]\label{lem:restriction}
固定上下文与切片外的表。设 $\widetilde f(\widetilde L,F_{\mathcal D})$
是 $\widetilde H^{2s}$ 与切片表上的函数。设其平移傅里叶支撑具有
公共对和性质
\[
 \widetilde Q_{g0}+\widetilde Q_{g1}
 =\widetilde Q_{h0}+\widetilde Q_{h1}
 \quad\text{对所有 }g,h\in[s].
\]
令 $f(L,F_{\mathcal D})=
\widetilde f(\iota L,F_{\mathcal D})$，其中 $\iota$ 作用于每个
平移。对每个标记 $q\in\mathcal A$，
\begin{equation}\label{eq:restriction-mass}
 m_q(f)\leq 2^{2s}
       \sum_{\substack{\widetilde q\in\widetilde{\mathcal A}\\
                       \rho(\widetilde q)=q}}
          \widetilde m_{\widetilde q}(\widetilde f).
\end{equation}
两边的质量都用 Equation~\eqref{eq:mass} 的公共对和定义，配各自的
仿射行空间。
\end{lemma}

\begin{proof}
对每组固定的切片表，傅里叶反演与特征正交性给出
\[
 \widehat f(Q)=
 \sum_{\substack{\widetilde Q\\
                  \iota^*\widetilde Q_{gb}=Q_{gb}\ \forall g,b}}
       \widehat{\widetilde f}(\widetilde Q).
\]
该和含 $2^{2s}$ 个指标。故其绝对值平方至多为各项绝对值平方和的
$2^{2s}$ 倍。对所有对和为 $q$ 的 $Q$ 求和：右侧每个非零系数的公共
对和 $\widetilde q$ 限制后为 $q$。此外，
\[
 \widetilde q(\one)=\widetilde q(\iota(\one))=q(\one)=1,
\]
故 $\widetilde q$ 是求值泛函，从而是 $\widetilde{\mathcal A}$ 中的
标记，且 $\rho(\widetilde q)=q$。每个提升指标只被计一次。对切片表
取平均即证 Equation~\eqref{eq:restriction-mass}。
\end{proof}

\subsection{端点染色者}

Equation~\eqref{eq:planted-score} 是对完整边元组的平均。为把它转
换成原博弈的染色，我们把输入边放在被植入的坐标处，并公开所有其他
边。零系数保证所需的上下文提升不需要输入边上的投影。

\begin{proposition}[解码]\label{prop:decoding}
设 Proposition~\ref{prop:extraction} 的抽取结论以常数 $p_*,\theta>0$
成立。若 $T$ 满足 Equation~\eqref{eq:decoder-repetition}，则原仿射
博弈的值至少为
\begin{equation}\label{eq:decoded-value}
 \Val(\mathcal G)\geq 2^{-2s}\frac{p_*\theta}{2}.
\end{equation}
\end{proposition}

\begin{proof}
我们构造只使用自身端点与公共随机性的随机化染色者。整个博弈与固定
的割表都作为背景数据可用。为给出博弈值的下界，不需要对染色者的
效率作任何断言。

\paragraph{共享数据与 $Y$ 染色者。}
给定来自原边律的随机输入边 $e=(x,y)$，选均匀的公开标记
$a_*\in[T]$。在每个其他槽位，公开采样一条独立边（含端点与投影）。
把输入边插入被标记的槽位。其投影只在分析中使用；不向任何染色者
透露。在任意标记的条件下，所得完整边元组具有乘积边律；故该元组
与标记独立。每个染色者知道自己的端点元组，且两者都知道被标记槽位
之外的所有边。

按上述植入实验公开采样深度、切片、节点、切片外的表、上下文、潜在
行与掩码。行采样使用固定的字母表维数、公开的树坐标与标记，且与
输入边独立。连同边元组，这些数据恰具有 Equation~\eqref{eq:planted-score}
使用的植入律。$Y$ 染色者形成其端点元组 $\underline y$，由
Equation~\eqref{eq:decoder-choice} 计算标记 $q^*$，并输出其被标记
分量 $q^*_{a_*}$。该计算用到的所有量对此染色者可用：特别是定义
$\bar G_{\underline y}$ 的条件平均可以由作为固定背景数据的博弈与
割算出。它不使用实际的被标记边或其投影。

\paragraph{$X$ 染色者。}
为定义 $X$ 染色者，把 $\mathcal A$ 的被标记因子换成 $X$ 字母表：
\[
 \widetilde{\mathcal A}
 =\prod_{a\ne a_*}\F^n\ \times\ \F^{n+1},
\]
因子按原槽位顺序排列。令 $\mathcal B=(\F^{n+1})^T$。完整投影分解为
\[
 \mathcal B\xrightarrow{\ \kappa\ }
 \widetilde{\mathcal A}\xrightarrow{\ \rho\ }\mathcal A.
\]
这里 $\kappa$ 在每个未标记槽位是公开边投影、在被标记槽位是恒等；
$\rho$ 在每个未标记槽位是恒等、在被标记槽位是未知映射 $\pi_e$。
于是 $\kappa$ 为 $X$ 染色者所知，而 $\rho$ 只在分析中使用。

植入的零意味着 $\Gamma$ 与 $R^{\mathrm{new}}$ 中的每行在被标记
槽位的线性系数块为零。$E$、$E_\Gamma$ 与 $\Gamma'=\Gamma+E_\Gamma$
亦然，因为噪声掩码与上下文变换只作标量乘与行加法。因此
$\widetilde\Gamma'=\rho^*\Gamma'$ 为 $X$ 染色者所知：把每个未标记
系数与常数系数保持不变、把所有被标记系数置零即得。该描述不需要
$\pi_e$，包括其仿射平移部分。

$X$ 染色者形成有界原始函数
\[
 \widetilde G(\widetilde R,F)
 =\sigma^o_{\underline x}
   \bigl(W_{\widetilde{\mathcal A}}(\widetilde R,F)
                \circ\kappa\bigr),
 \qquad \widetilde R\in\widetilde H^k.
\]
把其上下文固定为 $\widetilde\Gamma'$、其切片外的表固定为公开的
$F_{\mathrm{out}}$，在 $\widetilde H$ 中留下全部 $2s$ 个平移自由。
记其质量为 $\widetilde m_{\widetilde q}$。它们非负，且由 Parseval
与 $\abs{\widetilde G}\leq1$，其对标记之和至多 1。染色者按每个
标记至少 $\widetilde m_{\widetilde q}$ 的概率采样标记
$\widetilde q$，把剩余概率指派给固定标记。它用额外的独立公开均匀
随机数从该分布采样，并输出 $\widetilde q$ 的被标记分量。该采样不
使用 $q^*$。

\paragraph{在采样边上比较。}
我们验证把这些概率与分数比较所需的限制假设。对任意
$h\in\widetilde H$，把一对中两个子节点的生成子都平移 $h$，会使门
$v$ 的输出在每个标记 $\widetilde l$ 处乘以 $(-1)^{h(\widetilde l)}$。
在两个不同的对中施加该平移，该输出（从而整张词）不变。这些运算只
平移自由平移并保持上下文。因此特征正交性表明 $\widetilde G$ 的每
个平移傅里叶指标都有公共对和。该论证对每组固定的门表都成立。

词的生成与行的拉回交换（逐点于标记）。由于
$\pi_{\underline e}=\rho\circ\kappa$，把 $\widetilde G$ 的每个自由
平移限制到 $\rho^*H$ 上，恰得具有上下文 $\Gamma'$ 与同样表的
$G_{\underline e}$。映射 $\rho$ 的纤维为二元，故
Lemma~\ref{lem:restriction} 适用并给出
\[
 \sum_{\rho(\widetilde q)=q^*}\widetilde m_{\widetilde q}
 \geq 2^{-2s}
 m_{q^*}(G_{\underline e};\Gamma',F_{\mathrm{out}}).
\]
在全部植入数据的条件下，左边下界了采样到的完整标记投影到 $q^*$
的概率。该事件特别蕴含被标记的输入边被满足。取平均并用
Equation~\eqref{eq:planted-score}，成功概率至少
$2^{-2s}p_*\theta/2$。

最后固定所有共享随机性，包括用于采样 $X$ 标记的均匀数。同一组固定
数据在每个端点定义这两个策略，因此它们给出原博弈的确定性染色。
其关于共享随机性的平均具有刚才证明的下界，故某个固定选择达到它。
这证明 Equation~\eqref{eq:decoded-value}。
\end{proof}

解码损失依赖 $s,p_*,\theta$，从而依赖目标割间隙，但不依赖字母表
维数或重复长度。因此，在把博弈可靠性选得小于
Equation~\eqref{eq:decoded-value} 右端之后，可以先固定字母表维数，
再由 Equation~\eqref{eq:decoder-repetition} 选 $T$。

\section{完成难度证明}\label{sec:completion}

参数顺序为
\[
 (t,\eps,\mu)\longrightarrow(s,m,p_*,\theta)
 \longrightarrow b\longrightarrow n\longrightarrow T\longrightarrow a.
\]
抽取常数在字母表维数 $n$ 已知之前就固定所需的博弈可靠性 $b$。
重复次数随后选定以隐藏被植入坐标，博弈的完备性误差 $a$ 最后选取。

\begin{proof}[\Cref{thm:gap} 的证明]
按定理固定 $t$ 与 $\eps$。取 $\delta=\eps$。由于
$\eps<(t-B(t))/4<(1-B(t))/4$，这满足 $0<\delta<1-B(t)$。选有理数
$0<\mu<1/2$ 使 $t\mu\le\eps/4$。
以 $t,\delta,\mu$ 应用 \Cref{prop:extraction}，固定
$s,m,p_*,\theta$，从而固定 $k,d,r$。选有理数
\[
 0<b<2^{-2s}\frac{p_*\theta}{2}.
\]
由 \Cref{prop:affine-games}，这固定字母表维数 $n$，使博弈具有任意
小的完备性误差。选整数 $T$ 满足
\[
 T\ge\frac{2^{n(2k-2s)}}{(p_*\theta)^2}.
\]
最后选正有理数 $a<1-b$ 足够小使 $2Ta\le\eps/4$，并从 NP 难的博弈
间隙 $\Val(\mathcal G)\ge1-a$ 对 $\Val(\mathcal G)\le b$ 出发。

在 YES 情形，\Cref{lem:completeness} 给出割值至少
$(1+t)/2-\eps/2$，强于所需。在 NO 情形，假设存在值至少
$(1+B(t))/2+\eps/2$ 的割。由 \Cref{lem:cut-stability}，其奇平均的
期望切片稳定性至少为 $B(t)+\eps$。
\Cref{prop:extraction,prop:decoding} 随后将蕴含
$\Val(\mathcal G)\ge2^{-2s}p_*\theta/2>b$，矛盾。故每个割的值都
小于 $(1+B(t))/2+\eps/2$。

为完备起见：该测试是确定性的多项式规模图构造。顶点集有
\[
 |X|^T\,2^{|\mathcal B|}
 =|X|^T\,2^{2^{(n+1)T}}
\]
个顶点。除 $|X|$ 外的所有因子对固定的 $t,\eps$ 都是常数。端点与边
元组只有多项式多个，且所有词生成样本空间规模为常数。枚举每个测试
结果，把其概率加到相应边权上。初始博弈权重、$t$ 与 $r$ 都是有理
数，故这些权重具有多项式位复杂度。条件边概率同样具有多项式位复杂
度。所有字母表与词都可以在相对输入规模为常数的时间内枚举。因此所
得图可在确定性多项式时间内算出。它具有有理边分布，并满足所述的、
带 $\eps$ 损失的较弱界。
\end{proof}

\begin{proof}[\Cref{thm:main} 的证明]
把 $\rho=-t$ 代入 \eqref{eq:gw-constant} 给出
\begin{equation}\label{eq:gw-ratio}
 \GW=\min_{-1<t\le1}\frac{1+B(t)}{1+t},
\end{equation}
其中同一公式对负 $t$ 定义 $B(t)=2\arcsin(t)/\pi$。
当 $-1<t\le0$ 时比值至少为 1，而当 $0<t<1$ 时小于 1。在 $t=1$ 处
等于 1，当 $t\downarrow-1$ 时发散。故其最小值在 $(0,1)$ 内取得。

固定 $\GW<\alpha\le1$。由连续性与有理数的稠密性，选有理数
$t\in(0,1)$ 使
\[
 \alpha\frac{1+t}{2}>\frac{1+B(t)}{2}.
\]
选正有理数 $\eps$ 与 $\epsilon'$ 足够小，使 $\eps<(t-B(t))/4$ 且
\begin{equation}\label{eq:strict-final-gap}
 \alpha\left(\frac{1+t}{2}-\eps-\epsilon'\right)
 >\frac{1+B(t)}{2}+\eps+\epsilon'.
\end{equation}
应用 \Cref{thm:gap}。删除自环并保留原权重尺度，使剩余权重之和至多
为 1。以误差 $\epsilon'$ 应用 \Cref{lem:unweighted}，产生简单无权
图：在公共尺度 $Q$ 上，其最大割与原最大割权重相差至多 $\epsilon'$。
在该图上，$\alpha$-近似算法在 YES 情形返回缩放值至少为
\eqref{eq:strict-final-gap} 左端的割。在 NO 情形，没有割的缩放值
超过其右端。取严格介于二者之间的固定有理阈值即可区分 NP 难的承诺
情形。定理得证。
\end{proof}


\appendix
\section{由二对一博弈定理给出的仿射博弈}
\label{sec:affine-games}

Max-Cut 归约需要超出通常数值型二对一博弈间隙的两个特性：投影必须是
$\F$ 上的仿射映射，且字母表必须在选取完备性误差之前固定。
我们从 \cite[Sections~4--5]{DKKMS25} 的构造与可靠性分析推导这些
特性。其 Grassmann 一致性假设由 Khot、Minzer 与 Safra 的扩张定理
~\cite{KMS23} 以及 Barak、Kothari 与 Steurer 的扩张到一致性蕴含
~\cite[Theorems~9--11]{BKS19} 提供；\cite{DKKMS25} 发表版的引言
记录了这一补全。

我们先描述折叠博弈及其仿射投影。源的顶点定义排除了其显示采样器
可能产生的某些子空间。我们指明由此产生的条件化，补全其可靠性证明
所用的表，并给出带所需例外子空间估计的不变列表论证。对第二个
玩家，我们在商空间上使用有界列表，把每条限制到建议空间的部分分开
处理。完备性则通过最后选取初始 3Lin 误差得到。

\subsection{折叠博弈及其边律}

\paragraph{起始实例。}
\cite{DKKMS25} 的 Theorem~4.1（由 H\aa stad 定理~\cite{Hastad01}
得到）给出绝对常数 $s_*<1$：对每个足够小的 $\varepsilon>0$，区分
值至少 $1-\varepsilon$ 与值至多 $s_*$ 的正则二元 3Lin 实例是 NP 难
的。正则性包括变量的出现次数有界、每条方程三个不同变量、以及两个
不同方程至多共享一个变量。

现设 $\ell\ge2$、$k_0\ge\ell$ 与 $\beta\in(0,1)$（$\beta$ 有理）
为固定参数。两个字母表维数将是 $\ell$ 与 $\ell-1$，故
\Cref{prop:affine-games} 中的维数为 $n=\ell-1$。下文的可靠性论证
只用 $b$ 与绝对常数 $s_*$ 选取 $\ell,k_0,\beta$。

\paragraph{空间、标记与折叠。}
设 $\mathcal Z$ 为 3Lin 实例的变量集，令 $V_0=\F^{\mathcal Z}$，
把每条方程 $e$ 视为
\[
  \langle x_e,\alpha\rangle=b_e,
  \qquad x_e\in V_0,\quad b_e\in\F,
\]
其中 $x_e$ 的三个变量坐标为 1。\emph{合法}元组 $U$ 由 $k_0$ 条
互不相交的方程组成，并附加 \cite[Section~4.2]{DKKMS25} 施加的分离
条件：$U$ 中不同方程的变量不出现在任何输入方程中。令
\[
  V_U=\F^{\operatorname{var}(U)},
  \qquad H_U=\operatorname{span}\{x_e:e\in U\}.
\]
向量 $x_e$（$e\in U$）线性无关，故 $\dim V_U=3k_0$、$\dim H_U=k_0$，
且存在唯一线性函数 $h_U:H_U\to\F$ 使 $h_U(x_e)=b_e$。

折叠前，一个左问题是 $(U,L)$，其中
\[
  L\leq V_U,\qquad \dim L=\ell,\qquad L\cap H_U=\{0\}.
\]
对变量子集 $V\subseteq\mathcal Z$，记 $V_V=\F^V\le V_0$；$V$ 的
采样规则见下。右问题是 $(V,L')$，其中 $L'\leq V_V$、$\dim L'=\ell-1$。
在关联对上有 $L'\leq L$。两边的答案分别是 $L$ 与 $L'$ 上的线性
函数，约束是限制。选定基后，未折叠字母表即等同于 $\F^\ell$ 与
$\F^{\ell-1}$。

\cite[Section~4.2]{DKKMS25} 的折叠把 $(U,L)$ 与 $(U_1,L_1)$ 放入
同一类，当
\[
 L+H_U+H_{U_1}=L_1+H_U+H_{U_1}.
\]
其 Lemma~4.2 证明这是等价关系，并对每个类 $\mathcal C$ 给出固定的
$\ell$ 维空间 $R_{\mathcal C}$ 使
\begin{equation}
  R_{\mathcal C}+H_U=L+H_U
  \quad\text{对每个 }(U,L)\in\mathcal C.
  \label{eq:affine-complements}
\end{equation}
两边维数都是 $\ell+k_0$；因此也有 $R_{\mathcal C}\cap H_U=\{0\}$。
折叠答案是线性函数 $\sigma:R_{\mathcal C}\to\F$。其在 $(U,L)$ 上的
展开是唯一延拓
\begin{equation}
  \widetilde\sigma_U(r+h)=\sigma(r)+h_U(h),
  \qquad r\in R_{\mathcal C},\quad h\in H_U
  \label{eq:affine-unfolding}
\end{equation}
在 $L$ 上的限制。

为清晰起见，这里的仿射性是\emph{答案空间}之间映射的性质。对
$z\in L$，唯一地写 $z=A_Uz+B_Uz$（$A_Uz\in R_{\mathcal C}$、
$B_Uz\in H_U$）。由 \eqref{eq:affine-complements}，映射
$A_U:L\to R_{\mathcal C}$ 是线性同构。于是展开为
\[
  \sigma\longmapsto\sigma\circ A_U+h_U\circ B_U,
\]
这是从 $R_{\mathcal C}^*$ 到 $L^*$ 的仿射双射。与满射限制
$L^*\to(L')^*$ 复合，得到线性部分核为 1 维的满射仿射映射，每个
纤维恰有两个元素。在固定基之后即证得所需的仿射承诺。
\cite{DKKMS25} 的 Claim~4.4 表明：当若干未折叠边聚合为同一折叠边
时，这些投影彼此一致。等价地，这些边可以带各自的权重分开保留。

\paragraph{边采样的有效性。}
我们指明显示采样器省略的对病态样本的处理。独立采样 $k_0$ 条方程
产生非合法元组的概率为 $O(k_0^2/|\mathcal Z|)$（由有界出现次数）。
对固定的 $k_0$ 这可以任意小。必要时取输入实例的足够多个不相交
副本；这保持其值与正则性。

给定合法的 $U$，对每条方程独立地：右变量集 $V$ 以概率 $1-\beta$
包含全部三个变量，否则包含均匀选取的一个变量。若有 $f_0$ 条方程
贡献全部三个变量，则
\[
  \dim V_V=k_0+2f_0,
  \qquad \dim(V_V\cap H_U)=f_0.
\]
对均匀选取的 $(\ell-1)$-空间 $L'\leq V_V$，数非零向量得
\begin{align*}
  \Pr[L'\cap H_U\ne\{0\}]
  &\leq
  \frac{(2^{f_0}-1)(2^{\ell-1}-1)}{2^{k_0+2f_0}-1}
  \leq 2^{\ell-k_0}.
\end{align*}
若 $L'\cap H_U=\{0\}$，则包含 $L'$ 的均匀 $\ell$-空间
$L\leq V_U$ 可通过在 $V_U/L'$ 中选非零向量得到。它非平凡地交
$H_U$，当且仅当该向量属于 $(H_U+L')/L'$。故
\[
  \Pr[L\cap H_U\ne\{0\}\mid L']
  =\frac{2^{k_0}-1}{2^{3k_0-\ell+1}-1}
  \leq 2^{\ell-2k_0}.
\]
我们把采样条件化在合法性与 $L\cap H_U=\{0\}$ 上。被丢弃的概率
至多
\begin{equation}
  \kappa=O(k_0^2/|\mathcal Z|)
          +2^{\ell-k_0}+2^{\ell-2k_0}.
  \label{eq:affine-invalid-probability}
\end{equation}
固定 $\ell$ 后增大 $k_0$，两个几何项可以任意小。取其和至多 $1/4$，
并且必要时使用输入的不相交副本，把非合法性概率取至多 $1/4$。于是
从原独立方程采样器保留一个样本的总概率 $q_{\rm all}$ 至少为
$1/2$。我们把这一总保留概率与已在合法性条件化之后的归一化区分
开来。

设 $P$ 为对合法 $U$ 条件化后的该采样器，令 $E=\{L\cap H_U=\{0\}\}$
与 $\eta=2^{\ell-k_0}+2^{\ell-2k_0}$。我们的边律是
$R=P(\,\cdot\mid E)$，故 $\TV(P,R)=P(E^c)\leq\eta$。在源的有效边
集上定义次概率权重
\[
 w(e)=\sum_{\substack{\omega\in E\\\omega\mapsto e}}P(\omega),
 \qquad q_{\rm edge}=\sum_e w(e)=P(E)\geq1-\eta.
\]
于是 $R(e)=w(e)/q_{\rm edge}$。对任意折叠染色，其在 $R$ 下的接受
率是其原始有效接受率除以 $q_{\rm edge}$。原始有效接受使用上述律
$P$（对合法 $U$ 条件化但不条件化于 $E$），无效结果收益为零。下面
用已发表可靠性分析中补全的表来界定这个量。

\subsection{补全的表与其列表编码的方程}

可靠性证明在每个 $\ell$ 维子空间上使用表，尽管折叠染色只在有效
顶点处给出答案。我们先延拓这些答案，同时保持 $V_U$ 与 $V_V$ 上的
表之间的一致性。再证明从这些表解码出的线性函数仍满足 $U$ 中的
方程。

\paragraph{补全表。}
对与合法 $U$ 一起出现的右变量集 $V$，定义
\[
  H(V)=\operatorname{span}\{x_e:e\text{ 为输入方程且 }
                                     \operatorname{var}(e)\subseteq V\}.
\]
支撑在 $\operatorname{var}(U)$ 中的每条输入方程都属于 $U$：使用
两个元组块的方程违反合法性，而同一块内的不同方程会与该块的方程
共享三个变量。由于 $U$ 的方程支撑两两不交，这证明
\begin{equation}
  H_U\cap V_V=H(V).
  \label{eq:affine-validity-common-side-space}
\end{equation}
特别地，对 $L\subseteq V_V$，$(U,L)$ 的有效性与包含它的合法元组
$U$ 无关。

固定一个折叠染色。记 $F_U[L]$ 为其在每个有效 $(U,L)$ 处的展开
答案，并在 $L\cap H_U\ne\{0\}$ 时为零线性函数。对每个 $V$：当
$L\cap H(V)=\{0\}$ 时用任一包含它的 $U$ 定义 $F_V[L]$，否则取零。
对有效的 $L$，具有同一 $L$ 的所有顶点 $(U,L)$ 属于同一折叠类；其
展开答案由 \cite[Claim~4.4]{DKKMS25} 彼此一致。
Equation~\eqref{eq:affine-validity-common-side-space} 对无效的 $L$
给出同样的无关性。故补全的表满足
\begin{equation}
  F_U[L]=F_V[L]\qquad(L\subseteq V_V,\ \dim L=\ell).
  \label{eq:affine-validity-completed-consistency}
\end{equation}
辅助表 $F_V$ 定义在 $\ell$-空间上；它不同于博弈在 $(\ell-1)$-空间
上的右答案表。

补全的 $F_U$ 还有第二个性质。固定维数 $\ell+k_0$ 的
$K\supseteq H_U$。满足 $L+H_U=K$ 的所有有效 $L$ 属于同一折叠类。
公共折叠标记及其 $(H_U,h_U)$ 延拓给出一个线性函数 $f_K:K\to\F$：
$f_K|_{H_U}=h_U$ 且对每个这样的 $L$ 有 $F_U[L]=f_K|_L$。
该条件对无效表项不施加限制。

\paragraph{一致性与列表参数。}
表 $F$ 的 Grassmann 一致性测试先选均匀的 $(\ell-1)$-空间，再选
两个包含它的独立均匀 $\ell$-空间，当二者在公共 $(\ell-1)$-空间上
表项一致时接受。记其接受概率为 $\operatorname{agreement}(F)$。

固定可靠性参数 $\delta\in(0,1)$。一致性引理
\cite[Lemma~3.7]{DKKMS25} 在一致性阈值 $\delta^2/8$ 处给出常数
$q,r,C$ 与正函数 $\alpha(\ell)$（如其 Lemma~5.3 的证明）。取阈值
\[
  \tau_r=\min\{C/2,1/2\},\qquad
  \tau_i=10^{-9}\tau_{i+1}^{12}\quad(0\le i<r).
\]
于是 $0<\tau_i\le1/2$，且一致性引理给出的对的严格高于相关阈值。
再固定第二组列表阈值
\[
 \zeta_r=\tau_0/2,\qquad
 \zeta_i=10^{-9}\zeta_{i+1}^{12}\quad(0\le i<r).
\]
两个向量都在选取 $\ell$ 之前固定；$\ell$ 固定后 $\alpha(\ell)$ 为
正。我们选 $\ell$ 足够大，满足该一致性引理以及普通列表界
~\cite[Lemma~3.12]{DKKMS25}（对带阈值 $(\zeta_i)$ 的
$(\ell-q)$-空间）。还要求 $q\le\ell-1$ 与
\begin{equation}
  2^{q-\ell}\le\tau_0/4.
  \label{eq:affine-validity-ell-margin}
\end{equation}
重复参数 $k_0$ 随后选取，足够大使所引的环境维数要求在 $k_0$ 到
$3k_0$ 之间的每个维数及其对某个 $q$-空间的商上都成立。

对固定的 $q$-空间 $Q$，列表记录「限制在超过指定比例的含 $Q$ 的
$\ell$-空间上与表一致、余维数至多 $r$ 的空间上的线性函数」。只保留
无出现延拓的对。下一个引理表明：对上述补全的表，只要 $Q$ 避开其
张成，这些极大对就尊重原方程。

\begin{lemma}[可容建议空间上的不变列表断言]
\label{lem:affine-validity-invariance}
设 $X$ 维数为 $3k_0$，$H\le X$ 维数为 $k_0$，$h:H\to\F$ 线性。设表
$F$ 给每个 $\ell$-空间指派一个线性函数，且对每个
$(\ell+k_0)$-空间 $K\supseteq H$ 存在线性函数 $f_K:K\to\F$ 使
$f_K|_H=h$，且只要 $L+H=K$ 就有 $F[L]=f_K|_L$。固定满足
$Q\cap H=\{0\}$ 的 $q$-空间 $Q$，并令
\[
  e_{k_0}=2^{\ell+r-2k_0+1}.
\]
假设 $q\le\ell$、$3k_0-r>\ell$、\eqref{eq:affine-validity-ell-margin}
与 $e_{k_0}\le\tau_0/4$。对余维数 $i\le r$ 的 $Q\subseteq W\subseteq X$，
称 $(g,W)$ \emph{出现}，若 $g:W\to\F$ 线性且
\[
  \Pr_{Q\subseteq L\subseteq W,\ \dim L=\ell}
       [F[L]=g|_L]>\tau_i.
\]
若它没有到严格更大空间的出现延拓，则称其为极大的。每个极大出现对
满足 $H\subseteq W$ 且 $g|_H=h$。
\end{lemma}

\begin{proof}
令 $H'=H\cap W$。在 $W/Q$ 中，像 $(H'+Q)/Q$ 的维数为 $\dim H'$
（因为 $Q\cap H=\{0\}$）。数均匀 $(\ell-q)$-空间中的非零向量得
\begin{align*}
 d_W&:=\Pr_{Q\subseteq L\subseteq W}[L\cap H\ne\{0\}]\\
 &\le
 \frac{(2^{\dim H'}-1)(2^{\ell-q}-1)}{2^{\dim W-q}-1}
 \le 2^{\ell+\dim H'-\dim W+1}
 \le e_{k_0}.
\end{align*}
把 $W$ 换成任何余维数至多 $r$ 的包含空间，同一界成立。

把满足 $Q\subseteq L\subseteq W$ 的有效空间 $L$ 按 $Y=L+H'$ 分
类。在固定类内，表是 $Y$ 上某个尊重 $h|_{H'}$ 的线性函数的限制；
这由限制 $f_{L+H}$ 得到。若 $g|_{H'}\ne h|_{H'}$，这两个函数之差
在 $H'$ 上非零。$Y$ 中包含 $Q$ 的对 $H'$ 的补是 $(\ell-q)$ 维商
空间上线性映射的图像。其 $\ell-q$ 个独立提升选择必须各满足一个
非平凡线性方程，两个函数才会一致。故该类内的一致性至多
$2^{q-\ell}$（若两函数已在 $Q$ 上不一致则为零）。总一致性至多
\[
  d_W+(1-d_W)2^{q-\ell}
  \le e_{k_0}+2^{q-\ell}\le\tau_0/2<\tau_i,
\]
矛盾。故 $g|_{H'}=h|_{H'}$。

于是存在到 $\widetilde W=W+H$ 的唯一延拓 $\widetilde g$，尊重
$h$。在 $W$ 与 $\widetilde W$ 中含 $Q$ 的有效 $\ell$-空间的两种
均匀分布下，$L$ 在
\[
  W/(W\cap H)\simeq\widetilde W/H
\]
中的像在含 $Q$ 之像的 $\ell$-空间中均匀。事实上，每个这样的像在
$W$ 中有同样多的含 $Q$ 补，即 $2^{(\ell-q)\dim(H\cap W)}$，在
$\widetilde W$ 中为 $2^{(\ell-q)\dim H}$。此外，一致性只依赖该
像：$f_{L+H}$ 与 $\widetilde g$ 在 $H$ 上一致，故其差在一个补上
为零当且仅当在每个补上为零。因此两种环境空间中在有效空间上的
条件一致性相同。由于 $W$ 中的一致性至少为 $\tau_i$，
$\widetilde W$ 中的一致性至少为
\[
  (1-d_{\widetilde W})\frac{\tau_i-d_W}{1-d_W}
  \ge(1-e_{k_0})(\tau_i-e_{k_0})
  \ge\tau_i-2e_{k_0}\ge\tau_i/2.
\]
若 $H\not\subseteq W$，则 $\widetilde W$ 的余维数 $j$ 严格小于
$i$，且 $\tau_j\le\tau_{i-1}<\tau_i/2$。这给出一个出现延拓，与
极大性矛盾。
\end{proof}

取 $X=V_U$、$H=H_U$、$h=h_U$ 应用该引理可知：只要
$Q\cap H_U=\{0\}$，从 $F_U$ 解码的每个极大对都尊重 $U$ 的所有
方程。接下来估计该条件在源的可靠性证明所用分布下的概率。

\subsection{可靠性、完备性与参数顺序}

\paragraph{建议分布。}
\cite[Section~5.2]{DKKMS25} 的外层博弈先采样 $(U,V)$，再从 $V_V$
的 $q$-空间中均匀取 $Q$。它把 $(U,Q)$ 与 $(V,Q)$ 送给两个玩家，
要求分别在 $\operatorname{var}(U)$ 与 $V$ 上给出赋值。当这些赋值
在 $V$ 上一致且第一个满足 $U$ 的每条方程时接受。若 $f_0$ 条方程
贡献全部三个变量，则 $\dim V_V=k_0+2f_0$、$\dim H(V)=f_0$。因此
在该实际分布下，对 $U,V$ 一致地有
\begin{equation}
 \Pr[Q\cap H_U\ne\{0\}\mid U,V]
 \le\frac{(2^{f_0}-1)(2^q-1)}{2^{k_0+2f_0}-1}
 \le2^{q-k_0}.
 \label{eq:affine-validity-advice-loss}
\end{equation}
我们继续增大 $k_0$ 直到
\begin{equation}
 e_{k_0}\le\tau_0/4,
 \qquad 2^{q-k_0}\le\delta\alpha(\ell)/8.
 \label{eq:affine-validity-k-margins}
\end{equation}

\paragraph{从博弈接受到非空列表。}
设某折叠染色在 $P$（对合法 $U$ 条件化但不条件化于 $L$ 有效性的
采样器）下原始有效接受率大于 $\delta$。用补全的 $F_U[L]$ 对照
$(V,L')$ 上的原右答案只会提高该接受率：有效表项不变，而无效结果
此前的收益为零。因此我们可以把 \cite[Lemma~5.3, Claim~5.6, and
Equation~(5.1)]{DKKMS25} 的一致性/列表与覆盖论证用于补全的表。
在调用不变性之前，这些论证使用完全线性表、其一致性以及
\eqref{eq:affine-validity-completed-consistency}；现在它们在每个
采样到的子空间上都有定义。

对固定的 $U,Q$，记 $\mathcal L_{U,Q}$ 为 $V_U$ 中含 $Q$ 的
$\ell$-空间上的均匀律。$\mathcal L'_{U,Q}$ 先在外层问题 $(U,Q)$
的条件下采样 $V$，再取 $V_V$ 中含 $Q$ 的均匀 $\ell$-空间。这正是
\cite[Lemma~4.11]{DKKMS25} 的条件律。事实上，设 $N_q(d)$ 为
$\F^d$ 中 $q$-空间的个数，则对 $Q\subseteq V_V$，均匀
$\ell$-空间 $L\le V_V$ 满足
\[
 \Pr[Q\subseteq L\mid V]=\frac{N_q(\ell)}{N_q(\dim V_V)}.
\]
于是把源的 $V,L$ 采样器以 $Q\subseteq L$ 为条件，会给每个 $V$
赋以正比于 $\Pr[V\mid U]/N_q(\dim V_V)$ 的权重，与外层博弈完全
一致。

设 $E_1$ 为事件 $\operatorname{agreement}(F_U)>\delta^2/8$，$E_2$
为
\[
 \TV(\mathcal L_{U,Q},\mathcal L'_{U,Q})
 \le 2^{\ell+5}\sqrt\beta\,k_0^{1/4},
\]
$E_3$ 为 $F_U$ 的 $Q$-列表在阈值 $(\tau_i)$ 下非空。事件 $E_2$
正是源 Equation~(4.3) 的光滑性条件。对 $E_3$：只要 $E_1$ 成立，
源的 Lemma~3.7 对 $\alpha(\ell)$ 比例的均匀 $Q$ 给出一个出现对；
任一这样的对都有极大的出现延拓。这只用第一列表的存在性，不用其
规模界。所引论证给出这些事件的联合概率至少 $\delta\alpha(\ell)/4$。
与 $E_4=\{Q\cap H_U=\{0\}\}$ 相交后概率至少
$\delta\alpha(\ell)/8$（由
\eqref{eq:affine-validity-advice-loss}--\eqref{eq:affine-validity-k-margins}）。
事件 $E_4$ 只依赖 $U,Q$。因此在 $U,Q$ 固定后，它不改变 $V$ 的
后续条件律。

\paragraph{第二个玩家的商列表。}
对每个问题 $(V,Q)$ 与每个线性型 $\gamma:Q\to\F$，固定一个线性
延拓 $p_\gamma:V_V\to\F$。设 $\nu:V_V\to V_V/Q$ 为商映射。用
\[
 T_\gamma[L/Q]=
 \begin{cases}
  \overline{F_V[L]-p_\gamma|_L},&F_V[L]|_Q=\gamma,\\
  0,&F_V[L]|_Q\ne\gamma
 \end{cases}
\]
在 $V_V/Q$ 的 $(\ell-q)$-空间上定义完全表 $T_\gamma$。
此处 $Q\subseteq L\subseteq V_V$、$\dim L=\ell$，上横线表示
$L/Q$ 上诱导的线性型；第一种情形下它有定义，因为该差在 $Q$ 上
为零。

对该完全表，使用带阈值 $(\zeta_i)$ 的普通极大列表：$(g,S)$ 出现
当且仅当 $S\le V_V/Q$ 余维数 $i\le r$ 且 $T_\gamma$ 在 $S$ 中
超过 $\zeta_i$ 比例的 $(\ell-q)$-空间上与 $g$ 一致。
由~\cite[Lemma~3.12]{DKKMS25}，每个这样的列表规模至多
$M_0=M_0(\ell-q,r,\zeta_r)$，与 $k_0$ 无关。
第二个玩家选均匀的 $\gamma\in Q^*$，再从该列表均匀选一对 $(g,S)$。
它把 $g$ 提升为 $\nu^{-1}(S)$ 上的线性型 $g\circ\nu+p_\gamma$，
再均匀延拓到 $V_V$，并输出相应赋值。若其列表为空，则用任意赋值。
整个过程只用 $(V,Q)$ 与固定的表。

零表项可能在辅助表中产生额外的出现延拓。这不造成困难：下面的
证明只需要某个指定线性型的延拓，而不需要把这些辅助极大对等同于
$F_V$ 的极大对。

\paragraph{两个策略的成功。}
第一个玩家从 $F_U$ 的 $\tau$-列表中均匀选极大对 $(f,W)$，并把
$f$ 均匀延拓到 $V_U$（列表为空时用任意赋值）。这个有限列表不必
有与 $k_0$ 无关的规模界：成功估计对其每个元素都成立。第二个玩家
使用刚构造的商列表。

固定 $E_1\cap E_2\cap E_3\cap E_4$ 中的 $U,Q$ 与第一个玩家的
任意选择 $(f,W)$（$i=\operatorname{codim}_{V_U}W\le r$）。由
\Cref{lem:affine-validity-invariance}，$H_U\subseteq W$ 且
$f|_{H_U}=h_U$。于是第一个玩家选的每个延拓都满足 $U$ 的所有
方程。这是源的解码论证中不变列表断言的唯一用途。

\cite[Lemmas~4.11 and~5.5 and the proof of Lemma~5.3]{DKKMS25} 的
覆盖步骤（使用 \eqref{eq:affine-validity-completed-consistency}）
给出：对 $V$ 以至少 $0.3\tau_i\ge0.3\tau_0$ 的概率，
\[
 \operatorname{codim}_{V_V}(W\cap V_V)=i,
 \qquad
 \Pr_{Q\subseteq L\subseteq W\cap V_V}
       [F_V[L]=f|_L]>\tau_i/2.
\]
此处 $L$ 的维数为 $\ell$。令 $S_0=W\cap V_V$ 与 $\gamma=f|_Q$。
上式中的每个等式 $F_V[L]=f|_L$ 在 $T_\gamma$ 中保持。由于
$\zeta_i\le\tau_0/2\le\tau_i/2$，商对
\[
 \bigl(\overline{f|_{S_0}-p_\gamma|_{S_0}},\ S_0/Q\bigr)
\]
对 $T_\gamma$ 出现，从而在该辅助表中有极大的出现延拓。用同一个
$p_\gamma$ 提升任何这样的延拓，得到与 $f$ 在整个 $S_0$ 上一致的
函数。该结论只关乎函数延拓，不受零补全可能引入的额外一致性的
影响。

第二个玩家以至少 $1/M$ 的概率选中该标签与一个这样的极大延拓，
其中 $M:=2^qM_0$ 与 $k_0$ 无关。其最终答案与 $f$ 在
$W\cap V_V$ 上一致。第一个玩家的均匀延拓以概率
$2^{-i}\ge2^{-r}$ 与该答案在整个 $V_V$ 上一致：商
$V_V/(W\cap V_V)$ 的维数为 $i$。于是，在 $E_1\cap E_2\cap E_3\cap E_4$
的条件下，该策略获胜概率至少为
\[
 c_*:=2^{-r}\frac{0.3\tau_0}{M}>0.
\]
总成功概率至少为 $\delta\alpha(\ell)c_*/8$，这是与 $k_0$ 无关的
正常数。\cite[Lemma~5.4]{DKKMS25} 的上界在 $\beta k_0\to\infty$
时趋于零。取足够大的 $k_0$ 即得与源中相同的矛盾，证明原始有效
接受率至多为 $\delta$。

\paragraph{先定可靠性后定完备性。}
对 \Cref{prop:affine-games} 中的博弈，取 $\delta=b/4$ 并保留有效的
原始权重。由前述拒绝估计，其总质量 $q_{\rm edge}$ 至少为 $1/2$。
归一化后
\[
  \Val\le\delta/q_{\rm edge}\le b/2<b.
\]
上文的每个限制都只涉及 $b$ 与由它固定的参数。选定 $\ell$ 后，取
$k_0=j^4$、$\beta=j^{-3}$（$j$ 足够大）。则 $\beta k_0\to\infty$、
$\beta\sqrt{k_0}\to0$、$\sqrt\beta\,k_0^{1/4}\to0$。这些极限满足
外层博弈的可靠性与覆盖要求（包括 $2^\ell\beta\le1/8$）以及所有
显示的余量。特别地，字母表维数 $n=\ell-1$ 已只用 $b$ 固定。

\paragraph{字母表固定下的完备性。}
设 $\alpha:\mathcal Z\to\F$ 满足除 $\varepsilon$ 比例外的所有输入
方程，并把其线性延拓也记为 $\alpha:V_0\to\F$。给折叠左问题
$\mathcal C$ 答案 $\alpha|_{R_{\mathcal C}}$，给 $(V,L')$ 答案
$\alpha|_{L'}$。只要采样到的 $U$ 的每条方程都被满足，就有
$h_U=\alpha|_{H_U}$，故 \eqref{eq:affine-unfolding} 恰为
$\alpha|_{R_{\mathcal C}+H_U}$。因此该边被满足。此论证针对采样到
的边；不要求其折叠类的每个成员都涉及被满足的方程。

在任一条件化步骤之前，并界给出采样到未满足方程的概率至多
$k_0\varepsilon$。由于总保留概率 $q_{\rm all}$ 至少为 $1/2$，
条件化后的博弈值至少为
$1-k_0\varepsilon/q_{\rm all}\ge1-2k_0\varepsilon$。
在为可靠性固定 $\ell,k_0$ 之后，在初始 3Lin 间隙中取
$\varepsilon\leq a/(2k_0)$。这给出所述完备性而不改变 $n$。

\paragraph{有效性与规模。}
单个问题中的所有空间维数有界。输入方程的 $k_0$ 元组只有多项式
多个，且每个元组只有常数多个子空间、基与局部采样结果。折叠等价
关系可以由线性代数显式检验。为有效求出 $R_{\mathcal C}$，从类中
任一 $(U,L)$ 出发，枚举 $L+H_U$ 内的 $\ell$-空间（其维数为固定
常数 $\ell+k_0$），对每个类成员用 \eqref{eq:affine-complements}
检验每个候选。Lemma~4.2 保证有成功的候选。随后高斯消去计算基与
仿射边映射。有限采样分布、其条件化与平行边的聚合都具有多项式位长
的有理权重。这些观察给出命题所断言的确定性多项式归约。
\Cref{prop:affine-games} 证毕。

\section{转换到简单无权图}\label{sec:unweighted}

我们给出一个显式的确定性转换，使有理边权与自环不引入额外的复杂度
假设。比较使用原始的、未归一化的权重尺度。我们把每个顶点换成一颗
星团（cluster），并把每条带权边近似为一张二部图：其在每个矩形上的
边数以一致的绝对误差逼近加权密度的预测。任意割可以分裂每颗星团；
其星团比例于是描述原图的一个随机割，因此矩形估计控制放大图的所有
割，而不仅在星团上为常数的割。

\begin{lemma}\label{lem:unweighted}
设 $G$ 是 $N$ 个顶点的图，无序不同顶点对上有非负有理权重 $w_{uv}$，
且 $\sum_{u<v}w_{uv}\le1$。对每个固定 $\epsilon'>0$，可以在确定性
多项式时间内构造简单无权图 $G'$ 与整数尺度 $Q$，使
\[
 \left|\frac{\operatorname{MaxCut}(G')}{Q}
             -\operatorname{MaxCut}(G)\right|\le\epsilon'.
\]
当输入图有自环时，同样的结论在删除自环、且不对其余权重重新归一化
之后成立。
\end{lemma}

\begin{proof}
选素数 $p>\max\{2,2N^2/\epsilon'\}$。多项式规模的素数可以逐个检验
整数确定性地找到，用 Bertrand 公设界定搜索范围。把每个顶点换成以
$\mathbb F_p^6$ 索引的星团。对每个 $u<v$，设 $I_{uv}\subseteq\mathbb F_p$
为任意固定的、含 $\lfloor p w_{uv}\rfloor$ 个剩余类的集合，例如按
通常顺序取前面那些。当且仅当
\[
 x\cdot y\in I_{uv}
\]
时，把星团 $u$ 中的 $x$ 与星团 $v$ 中的 $y$ 相连。星团内部没有边。
这定义了 $Np^6$ 个顶点上的简单图，可通过枚举所有可能的对构造。
取 $Q=p^{12}$。

我们先估计它在两颗星团之间任意矩形上的边数。固定剩余类集 $I$ 与
子集 $S,T\subseteq\mathbb F_p^6$。记 $\psi(a)=\exp(2\pi i a/p)$，
傅里叶反演给出
\[
 \one_I(a)=\sum_{h\in\mathbb F_p}\beta_h\psi(ha),\qquad
 \beta_h=\frac1p\sum_{a\in I}\psi(-ha),\qquad
 |\beta_h|\le1.
\]
对 $h\ne0$，矩阵 $M_h(x,y)=\psi(hx\cdot y)$ 由特征正交性满足
\[
 M_hM_h^*=p^6 I,
\]
故其算子范数为 $p^3$。因此
\begin{align*}
 \left|\#\{(x,y)\in S\times T:x\cdot y\in I\}
                -\frac{|I|}{p}|S||T|\right|
 &\le\sum_{h\ne0}|\beta_h|p^3\sqrt{|S||T|}\\
 &\le p^{10}.
\end{align*}
取 $I=I_{uv}$，除以 $p^{12}$ 并用
$\bigl||I|/p-w_{uv}\bigr|\le1/p$，得
\begin{equation}\label{eq:rectangle-discrepancy}
 \left|\frac{e_{uv}(S,T)}{p^{12}}
       -w_{uv}\frac{|S|}{p^6}\frac{|T|}{p^6}\right|
 \le\frac1p+\frac1{p^2}.
\end{equation}

现在固定 $G'$ 的任意割，设 $a_u\in[0,1]$ 为星团 $u$ 在其正侧的
比例。星团 $u,v$ 之间的跨越边构成两个矩形。
Equation~\eqref{eq:rectangle-discrepancy} 于是把其缩放边数与
\[
 w_{uv}\bigl(a_u(1-a_v)+(1-a_u)a_v\bigr)
\]
比较，误差至多 $2/p+2/p^2$。对至多 $N(N-1)/2$ 个顶点对求和的代价
至多 $N(N-1)(1/p+1/p^2)$。上式显示的加权表达式对顶点对求和后，恰
是以概率 $a_u$ 独立地给顶点 $u$ 指派正号所得的期望割权重，因而至多
$\operatorname{MaxCut}(G)$。这证明了对 $G'$ 每个割的所需上界。

反之，把每颗星团中的所有顶点都取 $G$ 的最大割的符号。同一矩形界
把其缩放权重与 $\operatorname{MaxCut}(G)$ 比较到同样的误差，给出
反向不等式。$p$ 的选取保证
$N(N-1)(1/p+1/p^2)<\epsilon'$，断言得证。
自环对每个割的贡献为零，可在开头删除。
\end{proof}


\bibliographystyle{plain}
\bibliography{references}
\end{document}
